\documentclass[11pt,a4paper]{article}

\usepackage[T1]{fontenc}
\usepackage[utf8]{inputenc}
\usepackage{lmodern}
\usepackage{amsmath,amssymb,amsthm}
\usepackage{booktabs}
\usepackage{array}
\usepackage{microtype}
\usepackage{xcolor}
\usepackage[margin=1.05in]{geometry}
\usepackage[numbers,sort&compress]{natbib}
\usepackage[hidelinks,pdfusetitle]{hyperref}
\usepackage{url}
\urlstyle{same}
\emergencystretch=1.2em

\newif\iftodo\todofalse
\newcommand{\todo}[1]{\iftodo{\color{red!80!black}\sffamily\bfseries [TODO: #1]}\fi}

\theoremstyle{plain}
\newtheorem{theorem}{Theorem}[section]
\newtheorem{lemma}[theorem]{Lemma}
\newtheorem{proposition}[theorem]{Proposition}
\newtheorem{conjecture}[theorem]{Conjecture}
\newtheorem{corollary}[theorem]{Corollary}
\theoremstyle{definition}
\newtheorem{definition}[theorem]{Definition}
\theoremstyle{remark}
\newtheorem{remark}[theorem]{Remark}

\newcommand{\file}[1]{\texttt{#1}}
\newcommand{\Z}{\mathbb{Z}}
\newcommand{\PO}{P\!=\!O}
\DeclareMathOperator{\Placed}{Placed}

\title{Edge-graceful labellings of odd-order trees with few vertices of degree two}
\author{Lingsen Meng\\[2pt]
  \small Department of Earth, Planetary, and Space Sciences,\\[-2pt]
  \small University of California, Los Angeles, CA 90095, USA\\[-2pt]
  \small \texttt{lsmeng@g.ucla.edu}}
\date{6 September 2026}

% generated by scripts/paper_numbers.py -- do not edit
\newcommand{\AConst}{1.9\cdot 10^{8}}
\newcommand{\NumMaxWidth}{29}
\newcommand{\NumLatticeM}{24}
\newcommand{\NumParams}{7}
\newcommand{\NumCoefBound}{64}
\newcommand{\NumPhiMax}{10}
\newcommand{\NumPhiMin}{6}
\newcommand{\NumPlacedCoef}{26}
\newcommand{\NumPlacedCoefSink}{60}
\newcommand{\NumSinkBoxLin}{83{,}520}
\newcommand{\NumSinkBoxCon}{21{,}600}
\newcommand{\NumSinkLabLin}{37{,}416{,}960}
\newcommand{\NumSinkLabCon}{9{,}676{,}800}
\newcommand{\NumDetMax}{144}
\newcommand{\NumAdjMax}{216}
\newcommand{\NumCoefBoundEff}{27{,}648}
\newcommand{\NumLatticeMEff}{3{,}456}
\newcommand{\NumCoordLin}{36{,}192}
\newcommand{\NumCoordCon}{21{,}600}
\newcommand{\NumPairLin}{15{,}634{,}944}
\newcommand{\NumPairCon}{26{,}161{,}920}
\newcommand{\NumPairs}{11{,}927}
\newcommand{\NumGateFail}{0}
\newcommand{\NumGateFailRealizable}{0}
\newcommand{\NumContextList}{9{,}427}
\newcommand{\NumStates}{5{,}398{,}800}
\newcommand{\NumEmittable}{2{,}790}
\newcommand{\NumEmittableWithFamily}{2{,}790}
\newcommand{\NumEmittableMissing}{0}
\newcommand{\NumCatContexts}{5{,}717}
\newcommand{\NumCatFamilies}{12{,}748}
\newcommand{\NumEmittableFamilies}{6{,}698}
\newcommand{\NumEsixContexts}{2{,}648}
\newcommand{\NumEsixDenom}{2{,}789}
\newcommand{\NumNotEsixContexts}{141}
\newcommand{\NumTokenContexts}{1{,}724}
\newcommand{\NumAttempted}{14{,}047}
\newcommand{\NumSolverRuns}{30{,}485}
\newcommand{\NumReceived}{5{,}446}
\newcommand{\NumVariantContexts}{1{,}411}
\newcommand{\NumVariantFamilies}{2{,}588}
\newcommand{\NumFreeChildContexts}{544}
\newcommand{\NumFreeChildCovered}{544}
\newcommand{\NumNeutralBlocks}{47}
\newcommand{\NumNeutralTargets}{135}
\newcommand{\NumNeutralBuilt}{128}
\newcommand{\NumNeutralVariantBuilt}{20}
\newcommand{\NumNeutralTokenBuilt}{8}


\begin{document}
\maketitle

\begin{abstract}
A tree $T$ of odd order $n$ is edge-graceful if its edges can be labelled
bijectively by the nonzero residues modulo $n$ so that the vertex sums are all
the residues modulo $n$.  Lee conjectured in 1989 that every tree of odd order
is edge-graceful.  The case of trees with at most one vertex of degree two
follows from the zero-sum partition theorem of Kaplan, Lev and Roditty, and
we recently settled the case of at most two such vertices.  We prove that
there are absolute constants $C$, $A$ and $n_0$ such that every tree of odd
order $n\ge n_0$ with $D$ vertices of degree two is edge-graceful whenever
$n>C(D+1)(\log n)^{A}$; in particular, for every fixed $k$, every large enough
tree of odd order with at most $k$ vertices of degree two is edge-graceful.
The proof combines a finite table of local labelling identities, found by an
exact constraint search and re-verified by an independent checker, with a
bottom-up scheduling argument and with the zero-sum matching theorem of
M\"uyesser and Pokrovskiy.  With the linear deletion bound for cyclic groups
that its authors state in their concluding section as a consequence of their
proof, the threshold becomes linear, $n>C(D+1)$.  The constants are not
effective.  We also reduce the linear
threshold to an elementary statement about permutations of an interval, which
we prove in several cases and verify exactly up to order $30$, and which would
make the constant explicit.
\end{abstract}

\section{Introduction}\label{sec:intro}

Let $T$ be a tree with $n$ vertices, $n$ odd.  An \emph{edge-graceful}
labelling of $T$ is a bijection $\lambda\colon E(T)\to\Z_n\setminus\{0\}$
such that the induced vertex sums $s(v)=\sum_{e\ni v}\lambda(e)$, taken modulo
$n$, form a bijection $V(T)\to\Z_n$.  Lo~\cite{Lo1985} introduced the notion
and Lee~\cite{Lee1989} conjectured that every tree of odd order is
edge-graceful.  The parity condition is necessary, because the sum of all
edge labels is $\binom{n}{2}$, which is $\equiv 0$ modulo $n$ only for odd
$n$.  The conjecture is surveyed
in Gallian's dynamic survey~\cite{GallianDS6}.  The strongest general result
we are aware of is due to Kaplan, Lev and Roditty~\cite{KaplanLevRoditty2009},
who proved that $\Z_n\setminus\{0\}$, $n$ odd, admits a zero-sum partition into
parts of any prescribed sizes at least two, and deduced that every odd-order
tree with at most one vertex of degree two is edge-graceful.  Their argument
identifies the obstruction exactly, since a vertex of degree two would require
a zero-sum block of size one, which is impossible.  In \cite{Meng2026two} we
extended their theorem to trees with at most two vertices of degree two, by
combining zero-sum partitions with perfect and hooked Langford sequences and
a case analysis verified computationally for all odd $n\le 5001$.  The
argument there is specific to the two degree-two vertices and does not
extend to three.

In this paper we show that the number of degree-two vertices, rather than
their placement, is what limits the zero-sum approach, provided that it is
small compared with the order of the tree.

\begin{theorem}\label{thm:main}
There are absolute constants $C$, $A$ and $n_0$ such that every tree $T$ of
odd order $n\ge n_0$ with $D$ vertices of degree two is edge-graceful whenever
$n>C(D+1)(\log n)^{A}$.  In particular, for every $k\ge 0$ there is an $n_k$
such that every tree of odd order $n\ge n_k$ with at most $k$ vertices of
degree two is edge-graceful.
\end{theorem}

The cases $k=0$ and $k\le2$ of the second statement are the theorems of
\cite{KaplanLevRoditty2009} and \cite{Meng2026two}, respectively, with
$n_k=1$ there.  The constants come from an asymptotic matching theorem and
are not effective, and the theorem says nothing about trees in which the
degree-two vertices have positive density, e.g.\ caterpillars with many
subdivided edges, the regime of the full conjecture.

The only external input is a zero-sum matching theorem in $\Z_n$ that
tolerates a set of deleted residues, and it is the size of that set, not the
matching itself, that determines the threshold.  The printed statements of
M\"uyesser and Pokrovskiy~\cite{MuyesserPokrovskiy2025} tolerate
$n/(\log n)^{A}$ deletions and give Theorem~\ref{thm:main}.  For cyclic
groups, however, the authors state in their concluding Section~7 that a linear
number of deletions can be tolerated: the only source of a logarithm in their
proof is a bound on the length of a commutator product, which is empty in an
abelian group, and they record the consequence $g(n)=\Omega(n)$ for $\Z_n$
in their Problem~7.4.  With that bound our argument gives a linear threshold.

\begin{theorem}\label{thm:linear}
There are absolute constants $C$ and $n_0$ such that every tree $T$ of odd
order $n\ge n_0$ with $D$ vertices of degree two is edge-graceful whenever
$n>C(D+1)$.  In particular, for every $k\ge0$, every odd-order tree with at
most $k$ vertices of degree two and more than $C(k+1)$ vertices is
edge-graceful, since $D\le k$.
\end{theorem}

The proof of Theorem~\ref{thm:linear} uses the zero-sum matching lemmas of
\cite{MuyesserPokrovskiy2025} in the linear form for cyclic groups that the
authors state in their Section~7; that form is not a numbered theorem there,
and Remark~\ref{rem:mplinear} records exactly what is used and why it follows
from their proof.  The qualitative statement, that a bounded number of
degree-two vertices is no obstruction once the tree is large, does not depend
on it at all.  In Section~\ref{sec:linearroute} we reduce the linear
threshold, with an explicit constant and without any use of
\cite{MuyesserPokrovskiy2025}, to a conjecture on permutations of an interval
that we prove in several cases and verify exactly up to order $30$.

Our approach separates the problem into a local part and a global part.  Locally,
each maximal path of degree-two vertices, together with the branch vertex that
owns it, is labelled by an explicit family of integer forms whose vertex sums
permute its own labels (we call this the $\PO$ property).  Such families are
found by our exact search over all output permutations and constraint
programming, and every family used in the proof is re-verified by a program
that shares no code with the search.  Globally, the families are assembled
bottom-up so that the multiset of special labels is closed under negation up
to at most one zero-sum triple and has size and magnitude linear in $D$.  The
remaining labels are then partitioned into zero-sum blocks of the prescribed
sizes by the random Hall--Paige theorem of M\"uyesser and
Pokrovskiy~\cite{MuyesserPokrovskiy2025}.  Therefore the local part is a
finite computation (\NumCatFamilies{} certified families), whereas the
global part is a scheduling lemma, whose regression is an executable
scheduler that we tested on all trees with at most nine vertices and on
thousands of random trees.

The paper is organised as follows.  Section~\ref{sec:setup} reformulates
edge-gracefulness as a permutation identity and fixes the vocabulary.
Section~\ref{sec:decomp} decomposes the tree and proves the counting lemmas.
Section~\ref{sec:cells} describes the local families and the finite
certificate table.  Section~\ref{sec:realization} proves the sequential
realization lemma.  Section~\ref{sec:dense} applies the dense completion and
proves Theorems~\ref{thm:main} and~\ref{thm:linear}.
Section~\ref{sec:linearroute} records the elementary route to the linear
threshold.  Section~\ref{sec:verification} describes the computations and the
certificate format.

\section{Set-up}\label{sec:setup}

Throughout, $T$ is a tree of odd order $n$ rooted at a vertex $\rho$ of degree
at least three (if $T$ is a path the explicit labelling $1,2,\dots,n-1$ along
the path is edge-graceful, so this case is excluded).  For $v\ne\rho$ let
$e_v$ denote the parent edge of $v$ and put $\lambda(e_\rho)=0$.  A vertex is
\emph{stationary} if the labels of its child edges sum to zero, so that
$s(v)=\lambda(e_v)$.  Let $W$ be the set of non-stationary vertices together
with $\rho$ and let $S=\{\lambda(e_w):w\in W\setminus\{\rho\}\}$, the
\emph{special support}.

\begin{lemma}\label{lem:sigma}
$\lambda$ is edge-graceful if and only if the multiset $\{s(w):w\in W\}$ equals
$\{0\}\cup S$.
\end{lemma}

\begin{proof}
For a stationary vertex $v\ne\rho$ the sum $s(v)$ equals the label of $e_v$,
so the stationary vertices contribute exactly the labels of their parent edges.
The remaining labels are $\{0\}\cup S$ and they must be produced by $W$.
\end{proof}

Every vertex of degree two is in $W$, because a single nonzero child label
cannot sum to zero.  Therefore our construction chooses $W$ of size $O(D)$,
arranges the identity of Lemma~\ref{lem:sigma} cell by cell, and leaves every
other internal vertex a zero-sum block of size zero or at least two.  We work
with integer labels of magnitude $O(D)$ and reduce modulo $n$ at the end.

\paragraph{Core and arms.}
Suppressing the degree-two vertices of $T$ gives the core $Q$, a tree whose
vertices have degree one or at least three in $T$, and a core edge carries a
length $\ell$, the number of suppressed vertices.  We process $Q$ bottom-up.
A subdivided child edge of a core vertex $v$ is an \emph{arm} of $v$ when its
lower endpoint is not an owner (defined below), and the \emph{head chain} of
the lower endpoint otherwise.  A direct child edge whose lower endpoint is not
an owner is a \emph{port}.

\paragraph{Cells and $\PO$.}
A \emph{cell} at a core vertex $v$ consists of the labels of $v$'s parent edge
(or of its head chain), of its residual arms and of some recruited ports, and
possibly of the edge to one or two active children.  Its \emph{physical
multiset} $P$ is the multiset of these labels.  Its \emph{output multiset} $O$
consists of the sum at $v$, the adjacent sums along each arm and chain, the
terminal label of each arm, and the label of each port.  A cell is
\emph{self-permuting}, written $\PO$, when the two multisets coincide, and
for a root cell the constant $0$ is adjoined to $P$.  Substituting a $\PO$ cell for
a stationary child edge preserves the identity of Lemma~\ref{lem:sigma}
because the shared edge is counted once on either side.

\section{Decomposition and counting}\label{sec:decomp}

\subsection{Owner-free neutralization}

\begin{lemma}[neutral units]\label{lem:units}
Two arms whose lengths have the same parity and whose residues modulo six are
not $\{1,3\}$, and the four-arm multisets with residues $(1,1,1,3)$ and
$(1,3,3,3)$, admit labellings whose adjacent sums and terminals permute their
own labels, whose first labels sum to zero, and whose support is closed under
negation, and each has two free integer parameters.  Two arms of lengths one and
three admit no such labelling.
\end{lemma}

The thirteen families for the residue representatives are found by our
constraint search of Section~\ref{sec:cells} (a root row with pinned zero
output) and re-verified by the independent checker.  They extend to all
lengths in the residue classes by the six-edge insertion, which lengthens an
arm by six while adding three inverse pairs, and the impossibility for lengths
one and three is an exact enumeration.  We call these labelled arm sets
\emph{neutral units}.  (Pairs whose residues are $\{1,3\}$ but whose lengths
are not, e.g.\ $(1,9)$ or $(3,7)$, also admit such labellings by the same
search.  We do not use them, since a pair of this type is carried instead in
the residual, which has families of its own.)

\begin{corollary}[six residual types]\label{cor:sixtypes}
The arm multiset of every core vertex can be partitioned into neutral units
and a \emph{residual} of one of six types, namely empty, one odd arm, one even
arm, a pair with residues $\{1,3\}$, one odd and one even arm, or such a pair
together with one even arm.
\end{corollary}

\begin{proof}
Pair the even arms arbitrarily, leaving at most one even arm.  Among the odd
arms, pair every arm of residue five with any other odd arm, then pair arms of
residue one with each other and arms of residue three with each other.  What
remains is at most one arm of residue one and at most one of residue three.
Thus the odd part of the residual is empty, a single odd arm, or a pair with
residues $\{1,3\}$, and the even part is empty or a single arm.
\end{proof}

\subsection{Owners and contexts}

\begin{definition}\label{def:owner}
A nonroot core vertex is an \emph{owner} if its residual is nonempty, or if
exactly one active child remains after the cancellations described below and
no port is available to absorb it.  A vertex with empty residual and no
remaining active child is stationary, and its parent edge, if subdivided, is an
arm of its parent.
\end{definition}

At a vertex with $k\ge 2$ active children (children which are owners) the
heads are free parameters of the children's cells, so they are prescribed to
sum to zero (a mirror pair for $k=2$, a free zero-sum triple for $k=3$, pairs
and at most one triple otherwise), and none remains.  A mirror pair is closed
under negation, but a cancelling triple is not, and its negatives are placed at
a sink in the ordinary part of the tree exactly as the negatives of a token
are; Appendix~\ref{app:realization} counts these triples together with the
tokens.  A single remaining active
child with head $x$ is absorbed by a stationary vertex with $q$ ports by a
helper port labelled $-x$ ($q=1$ or $q\ge 3$) or by two ports labelled $a$ and
$-x-a$ ($q=2$).  Two exceptional child types cannot be absorbed in this way,
because every $\PO$ family of the child contains the antipode of its own head:
the residue-one letter and the $(1;\text{one port})$ bottom, and such children
are joined into their parent's cell.

The \emph{context} of an owner is the tuple $(h,R,p,c)$ where $h$ is the
length of its head chain reduced modulo six, $R$ its residual arm multiset
with lengths reduced modulo six, $p\in\{0,1,2,3\}$ the number of recruited
ports (never leaving exactly one unrecruited), and $c$ the child type, one of
\emph{none}, \emph{free}, the two joined types, or the two-input type used
when two active heads are kept.  The root has context $(\mathrm{root},R,p,c)$.

Reduction modulo six is justified by the following identity, which lengthens
one arm of a cell by six and changes nothing else.

\begin{lemma}[six-edge insertion]\label{lem:insertion}
Let an arm of a cell carry the labels $a_0,\dots,a_\ell$, $a_0$ at the owner,
so that it contributes the outputs $a_0+a_1,\dots,a_{\ell-1}+a_\ell$ and
$a_\ell$.  Write $c=a_0$ and $d=a_1$ and insert the six labels
\[
  \tfrac{-2c-d}{3},\quad \tfrac{c-d}{3},\quad \tfrac{c+2d}{3},\quad
  \tfrac{2c+d}{3},\quad \tfrac{-c+d}{3},\quad \tfrac{-c-2d}{3}
\]
between $a_0$ and $a_1$.  They form three pairs of opposite forms, and the
seven outputs they create are exactly those six labels together with the
output $c+d$ that the arm had before.  For a port, an arm of length zero whose
only label is $c$ and whose output is $c$ itself, the same holds with the six
labels
\[
  \tfrac{c}{3},\quad \tfrac{-2c}{3},\quad \tfrac{4c}{3},\quad
  \tfrac{-c}{3},\quad \tfrac{2c}{3},\quad \tfrac{-4c}{3}
\]
and with $c$ in place of $c+d$.
\end{lemma}

\begin{proof}
For the first list the outputs are, in order,
$c+\frac{-2c-d}{3}=\frac{c-d}{3}$, $\frac{-2c-d}{3}+\frac{c-d}{3}=\frac{-c-2d}{3}$,
$\frac{c-d}{3}+\frac{c+2d}{3}=\frac{2c+d}{3}$,
$\frac{c+2d}{3}+\frac{2c+d}{3}=c+d$,
$\frac{2c+d}{3}+\frac{-c+d}{3}=\frac{c+2d}{3}$,
$\frac{-c+d}{3}+\frac{-c-2d}{3}=\frac{-2c-d}{3}$ and
$\frac{-c-2d}{3}+d=\frac{-c+d}{3}$, which is the asserted multiset, and the
three pairs are the first with the fourth, the second with the fifth and the
third with the sixth.  The port list is checked the same way.
\end{proof}

Thus a family for a context is also a family for the context with any one
arm six longer, and in particular a port may be replaced by an arm of length
six.  The new labels are forms in the two labels the insertion sits between,
so the identity is one statement about all cells at once.  The
insertion may make two labels of the cell equal for special values of $c$ and
$d$, which the search rules out by requiring the label forms to stay pairwise
distinct.  A twelve-label endpoint preserving insertion does the same for head
chains.

A second identity removes a much larger part of the search.  Most of the
contexts of the list differ from a narrower one only by a set of arms and
ports that can be made to pay for itself, and such a set may be adjoined to
any family at all.

\begin{definition}\label{def:block}
A \emph{block} is a multiset of arm lengths $\ell_1,\dots,\ell_k$ together
with a number $m\ge0$ of ports, and it carries $N=m+\sum_i(\ell_i+1)$ labels.
The block is \emph{neutral} if there are integers $c_1,\dots,c_N$ such that,
when its labels are taken to be $c_1t,\dots,c_Nt$ in a single parameter $t$,
\begin{enumerate}
\item[(N1)] the labels it contributes to the owner, namely its $m$ port labels
  and the first label of each of its arms, sum to zero;
\item[(N2)] the multiset of its $N$ labels equals the multiset of the $N$
  outputs at the vertices it introduces, that is the $m$ port leaves and, for
  an arm $a_0,\dots,a_\ell$, the outputs $a_0+a_1,\dots,a_{\ell-1}+a_\ell$ and
  $a_\ell$;
\item[(N3)] the multiset $\{c_1,\dots,c_N\}$ is closed under negation; and
\item[(N4)] the $c_i$ are nonzero and pairwise distinct.
\end{enumerate}
It is a \emph{token block} if it satisfies (N1), (N2) and (N4) and, instead of
(N3), three of its labels sum to zero and span a plane while the remaining
$N-3$ are closed under negation.  A token block therefore needs at least two
parameters, and $N$ is odd for it and even for a neutral block.
\end{definition}

Which of the two a given set of arms and ports can be is decided by parity
alone.  Two arms of lengths $a$ and $b$ carry $a+b+2$ labels, so a same-parity
pair can only be a neutral block and a pair of opposite parity can only be a
token block, which is the arithmetic behind the neutral units of
Section~\ref{sec:decomp}: an $\{o,e\}$ residual is exactly the residual that
forces a token.

\begin{lemma}[neutral adjunction]\label{lem:adjunction}
Let $B$ be a neutral block.  Then every family for a context $C$ is also a
family for the context obtained by adjoining the arms and the ports of $B$ to
the owner of $C$, with the same head, the same pinned input and the same
token.  If $B$ is a token block the same holds provided $C$ has at most one
token, and the enlarged family carries the triple of $B$ as its token when $C$
has none and as its second token when $C$ has one.
\end{lemma}

\begin{proof}
Give $B$ a parameter $t$ that the family for $C$ does not use.  Every label of
$C$ is then zero in $t$ and, by (N4), no label of $B$ is, so the labels of the
enlarged cell are nonzero, they are pairwise distinct across the two groups
because they differ in the coordinate $t$, and they are pairwise distinct
inside $B$ again by (N4).  By (N1) the owner keeps the output it had, and the
only other outputs that change are those created inside $B$, so $P=O$ for the
enlarged cell follows from $P=O$ for $C$ together with (N2).  Condition (N3)
splits the new labels into pairs of opposite forms, which extends the
signature partition by mirror pairs and leaves the token where it was.  For a
token block the same partition is extended by the mirror pairs together with
the new triple, which sums to zero and spans a plane as the signature demands,
and which is disjoint from the token of $C$ because it lives in the fresh
parameters.
Finally the head chain, the pinned input, the token and the ports of $C$ are
untouched, so the free-head condition, the input condition, the token rank
and each of the variant conditions of Section~\ref{sec:cells} hold for the
enlarged family exactly as they held before.
\end{proof}

The two blocks we use most often are neutral for a reason one can write down.

\begin{lemma}\label{lem:neutralpairs}
Two arms of the same length $\ell$ form a neutral block, and so do two ports.
In addition, two arms of lengths $a<b\le8$ of the same parity form a neutral
block for every such pair except $\{1,3\}$, which is not neutral; one port
together with an arm of even length $\ell\in\{4,6,8\}$ is neutral; one port
together with two arms of opposite parity is neutral; and two ports together
with an arm of odd length $\ell\in\{5,7\}$ are neutral.  In addition, two arms
of opposite parity $\{o,e\}$ form a token block for every pair with
$o\in\{3,5\}$, $e\in\{2,4,6\}$ and for $\{1,6\}$ and $\{2,5\}$, although
$\{1,2\}$ and $\{1,4\}$ do not.
\end{lemma}

\begin{proof}
For two arms of length $\ell$ take $a_i=(-2)^it$ on one and $a_i=-(-2)^it$ on
the other.  The outputs of the first are $a_i+a_{i+1}=-a_i$ for $i<\ell$
together with $a_\ell$, and those of the second are their negatives, so the
$2\ell+2$ outputs are exactly the $2\ell+2$ labels $\pm(-2)^it$,
$0\le i\le\ell$.  Condition (N1) holds because the two first labels are $t$
and $-t$, condition (N3) because the labels come in opposite pairs, and (N4)
because the $2^i$ are distinct.  Two ports labelled $t$ and $-t$ satisfy the
four conditions directly.  For the remaining blocks, and for the assertion
that $\{1,3\}$ admits no neutral labelling, the four conditions are a finite
integer feasibility question once a bound on the $c_i$ is fixed, and we
settled it by an exhaustive search over integer labels with coefficients at
most $8$ in absolute value; the resulting labels are listed in
Table~\ref{tab:blocks}, and each is re-verified from its labels alone, with no
reference to how it was found.
\end{proof}

That $\{1,3\}$ is the one exception does not appear to be an accident of the
search.  It stays infeasible when the block is allowed three parameters
instead of one and coefficients up to $14$ in absolute value.  We have not
proved that no labelling whatsoever works, and we do not claim it; the pair is
in any case the one our construction singles out everywhere else as well.  A $\{1,3\}$ block can still be carried
by a cell, as the catalogue shows; what fails is only the uniform statement
that it may be adjoined to an arbitrary family.

Lemma~\ref{lem:adjunction} is what keeps the catalogue within reach.  A
context with a retained neutral unit, or with a pair of recruited ports, is
almost always a narrower context with a block adjoined, and the wide contexts
are exactly the ones a solver finds hard.  Of the \NumNeutralTargets{}
contexts our searches had not reached, \NumNeutralBuilt{} are obtained in this
way from a narrower context that already had a family, at no search cost at
all, a further \NumNeutralVariantBuilt{} input-in-token variants are obtained
the same way, since by Lemma~\ref{lem:adjunction} the adjunction does not move
the pinned input out of the token, and the last \NumNeutralTokenBuilt{}
contexts, all of them with an $\{o,e\}$ residual and no port, are obtained from
a token block.  Together with the searches this leaves no context of
Lemma~\ref{lem:contexts} without a family.  The construction is
not trusted: the families it produces are written in the format the solver
emits and are re-verified by the independent checker of
Section~\ref{sec:verification} together with everything else.

\begin{table}[t]
\centering
\caption{The neutral blocks of Lemma~\ref{lem:neutralpairs}.  The labels are
read as multiples of a fresh parameter $t$, the ports first and then the arms
in the order listed, each arm from the owner outwards.  The equal pairs are
the closed form $\pm(-2)^i$; the others come from an exhaustive search and are
re-verified from the labels alone.}
\label{tab:blocks}
\small
\input{blocks_table}
\end{table}

\begin{lemma}[context completeness]\label{lem:contexts}
Every owner receives exactly one context from a finite list, and every
vertex of degree two lies in exactly one arm, head chain or neutral unit.
Exactly \NumEmittable{} of the contexts of the list occur.
\end{lemma}

The case analysis is given in Appendix~\ref{app:contexts}.  Its regression,
our executable scheduler, produces a context in the list for every owner of
every tree with at most nine vertices ($4{,}782{,}969$ labelled trees) and of
$169{,}000$ uniform random labelled trees, $20{,}000$ at each odd order
$n\in\{11,13,15,17,19,21,31,51\}$ and $3{,}000$ at each of
$n\in\{101,201,401\}$.  No owner of any of them is left without a context and
no context outside the list is produced.  That the root is a branch vertex is used, and not
only for convenience.  A vertex of degree two taken as the root can be left
with a residual that its arms and its one port cannot cover and with no active
child to spend on it, which is the configuration
$(h,\{1\},1,\mathrm{none})$ at the root; a branch vertex has a third core child
and cannot be in that position.  On $2{,}000$ random trees of orders $11$ to
$51$, rooting at a vertex of degree two left an owner without a context for
$196$ of them.  Rooting at a branch vertex never did: over $3{,}000$ random
trees of orders $11$ to $51$ we ran the decomposition from \emph{every} branch
vertex in turn, $23{,}000$ roots in all, and every one of them gave every owner
a context.  The choice among branch vertices is therefore free, as
Lemma~\ref{lem:contexts} states it, and not something the decomposition has to
search over.
In addition, the decision at a vertex depends only on its head-chain
length, its arm lengths, its number of ports and the types of its active
children.  We therefore enumerated all $\NumStates$ abstract local states with up to
three arms of lengths at most eight (and larger arm sets of short arms), up to
three ports, up to four active children of each of the possible types and
head chains of length at most six.  Each state was realised by a concrete tree
fragment, and we recorded the context emitted at the vertex.  The
\NumEmittable{} distinct contexts obtained are all in the list, and they are
exactly the contexts for which families are required in
Section~\ref{sec:cells}, whereas the remaining entries of the list are slack.

The cap on the number of active children is the one hypothesis of this
enumeration that the decomposition does not force, since an owner of a large
tree may carry arbitrarily many active children, and a tree on at most nine
vertices carries none with three.  We therefore ran the enumeration separately
at the caps two, three and four.  Raising the cap from two to three adds six
contexts, all of them roots with a single arm, and $416$ realizable ordered
pairs, whereas raising it from three to four adds nothing at all: the
$3{,}818{,}400$ states enumerated at the cap four emit exactly the same
contexts and exactly the same ordered pairs as the $1{,}204{,}800$ states at
the cap three.  Thus the emitted set is stationary between three and four,
which is what the analysis of Appendix~\ref{app:contexts} predicts, since the
root rule and the two-input rule reduce any number of active children to one
or two before a context is emitted.  We note that a stationary pair of caps is
evidence for, and not a proof of, stationarity at every cap; the proof is the
case analysis itself.

\begin{lemma}[charge]\label{lem:charge}
Let $q\le D$ be the number of subdivided core edges.  The number of owners is
at most $2q$, the special support satisfies $|S|\le 12D+3$, and the values
reserved for forced companions number at most $2\NumPhiMax q$.
\end{lemma}

\begin{proof}
Every owner is incident with a subdivided core edge, and such an edge has two
endpoints.  A subdivided edge of length $\ell$ carries $\ell+1$ labels, at
most $D+q$ in total.  Recruited ports contribute at most three per owner,
absorptions at most two per owner, and the final bridge triple three.  An owner
reserves one value for each ratio in $\Phi(F)$ of the family it is given, so
the number reserved is at most $\NumPhiMax$ per owner and at most
$2\NumPhiMax q$ in total, where $\NumPhiMax$ is the largest number of forced
ratios of any family in the catalogue (a choice rule that preferred families
with few ratios would lower it to $\NumPhiMin$, the largest over the contexts
of the smallest number in the menu, but our constructor does not use such a
rule and we do not rely on one).
\end{proof}

\section{Local families}\label{sec:cells}

\subsection{Symbolic families}

A \emph{family} for a context is a list of labels given as integer linear
forms, divided by a common denominator $L$, in parameters $y$ (the head),
$x$ or $x_1,x_2$ (inputs), token coordinates and internal parameters, such
that
\begin{enumerate}
\item[(E1)] $\PO$ holds identically in the parameters;
\item[(E2)] the labels other than the head split into pairs of opposite forms
  and at most two \emph{tokens}, triples of forms summing to zero (a second
  token is needed only for two-input cells whose non-head label count is
  even);
\item[(E3)] the label forms are pairwise distinct and nonzero;
\item[(E4)] the three coefficient vectors of each token span a plane, and a
  \emph{carrier} of at most two parameters on which they already span a
  plane is recorded;
\item[(E5)] the head is an independent parameter (nonroot contexts);
\item[(E6)] the only label depending on an input alone is the antipode of that
  input, and no non-head label is a rational multiple of the head.
\end{enumerate}
We call a family \emph{forced-antipode} when some non-head label is
identically minus the head.  Such a label necessarily lies in a token
$\{-x,a,x-a\}$ of the child, since its mirror is the head and the head is
excluded from the signature, and a parent can then absorb the head $x$
neither by a helper port nor by a mirror sibling, because the antipode of the
head is already used inside the child.  Two contexts of our catalogue carry
only forced-antipode families, the $(1;\text{one port})$ bottom and the
two-input cell with head chain of length two and residual $\{1\}$, and only
the second of them occurs, the first being folded into its parent.  Every
other context has a family that is not forced-antipode, so a single child
context is responsible for the whole of what follows.

A parent absorbs the head of a forced-antipode child through a token of its
own that contains the shared edge, $\{x,b,-x-b\}$, so that the two tokens are
negatives of each other once the child's carrier is matched, $a=-b$.  A
stationary parent with at least two ports does this with its port block.  An
owner uses an \emph{input-in-token} family of its context, found by the
same search with the shared edge required to be a token member (and with a
second, free token when the number of non-head labels is even).  A parent
whose context carries no such family instead absorbs the child by joining its
rows into its own cell, the joined child type A2 of
Appendix~\ref{app:contexts}, and the joined cell then has a family that is
not forced-antipode, so that it is an ordinary child of the next parent and
the absorption stops after one step.  Although the two rules look
interchangeable, the second is the one that is always available, and the
first is what keeps the cells small.

One parent context deserves a separate remark, since it shows that the
joining rule cannot be dispensed with.  Let $v$ be an owner with head chain
of length $\equiv 0$, one arm of length $1$, no ports and two kept inputs
$x_1,x_2$, write $a_1,a_2$ for its arm labels and $h$ for its head, and read
off the three outputs $h+a_1+x_1+x_2$, $a_1+a_2$ and $a_2$.  The leaf output
is $a_2$, and $a_1+a_2=a_1$ would force $a_2=0$, so $a_1+a_2=h$ and therefore
$h=-(x_1+x_2)$, which leaves $a_1=t$, $a_2=-(x_1+x_2)-t$ with $t$ free.  The
four non-head labels admit no token, and pairing $t$ with $-(x_1+x_2)-t$, or
$x_1$ with $x_2$, forces $h=0$.  Thus $t=-x_1$ or $t=-x_2$, and in
either case the labels are $-x_1,-x_2,x_1,x_2$.  Thus this owner always
carries the antipodes of both of its inputs, it has no input-in-token family,
and a forced-antipode child under it must be joined.  In the joined cell the
second child is pinned by one label whose output is that label itself, which
is exactly the role of a port, so the joint is searched as the context with
one more port and one active child.  Property (E6) is what makes adjacent cells compatible without any
case analysis, since a parent forces at most the antipode of its input and a
child never carries the antipode of its own head, so that no two forms of
adjacent cells coincide identically.  Property (E6) has two clauses and both
are needed for that argument, the first for the parent and the second for the
child.  We test them together and literally, over every family of every menu
and, for a two-input row, against both of its inputs; a separate program does
this, since our checker verifies (E1)--(E5) only.  By that test, for
\NumNotEsixContexts{} of the \NumEsixDenom{} contexts that occur and carry a
family our catalogue has no family with (E6), and for these the finitely many
adjacent pairs are checked individually by the gate of
Section~\ref{sec:verification}.  We note that the gate does not rely on this
count, since it enumerates every ordered pair of contexts and not only the
pairs that (E6) leaves open.

\subsection{The search and the checker}

Families are found by our constraint model over the integer coefficients of all
label forms: one Boolean per (output, label) pair encodes the permutation,
each true Boolean imposes the equality of two coefficient vectors, the
signature (mirror pairs, token, head, inputs, zero label) is encoded by further
Booleans, and every coefficient vector is required to be nonzero and distinct
from the others.  The model is solved with CP-SAT~\cite{ORTools}.  A solution
is a symbolic family, not a numerical labelling, and it is accepted only after
an independent program reconstructs the topology from the row specification,
recomputes the outputs, and re-verifies (E1)--(E5) with exact rational
arithmetic and at three random integer points.  Every family in the table
passes, and in addition the six-edge mother and the twelve-label insertion are
applied once to every arm and every head chain of every family and the result
is re-verified.  Table~\ref{tab:catalogue} summarises the catalogue.

\begin{table}[t]
\centering
\caption{The certificate table, by kind of context: the contexts that occur
(Lemma~\ref{lem:contexts}), those for which the catalogue carries a family,
the number of families, the contexts with a family satisfying (E6), and the
contexts whose chosen family carries a token.  The residue-two rigid row is a
constructor template and not a symbolic family, so it is counted as covered in
the second column and in none of the others.}
\label{tab:catalogue}
\begin{tabular}{lrrrrr}
\toprule
context kind & contexts & with a family & families & (E6) & token \\
\midrule
bottom cells & 954 & 954 & 1{,}563 & 913 & 493 \\
continuing cells, free child & 419 & 419 & 1{,}777 & 418 & 406 \\
continuing cells, joined child & 876 & 876 & 1{,}808 & 810 & 440 \\
two-input cells & 34 & 34 & 213 & 31 & 33 \\
root cells & 501 & 501 & 1{,}316 & 501 & 347 \\
\midrule
total & 2{,}784 & 2{,}784 & 6{,}677 & 2{,}673 & 1{,}719 \\
\bottomrule
\end{tabular}

\end{table}

Four exact impossibilities shape the alphabet.  A single arm of length four
with no ports has only one-parameter families (rays), so a residual even
singleton retains one neutral unit.  A head chain of length one or two with
ports only has no $\PO$ family at all, and a subdivided parent edge therefore
does not by itself make a vertex an owner.  The residue-one letter on a direct parent
edge has the unique family $(-x;\,-x-t,\,t;\,x)$, whose head is the antipode
of its input and is therefore not a free parameter, and it is for this reason
joined with its parent.  With a subdivided parent edge the same letter has
ordinary families with a free head, e.g.\ $(y;\,-y-x,\,-x,\,y+x;\,x)$ for
$h=1$, and it is then an owner of its own.  Finally, the two-input cell with
head chain of length two and residual $\{1\}$ has no family without the
antipode of its own head, over all denominators $L\le 12$, all numbers of
internal parameters up to three, all coefficient bounds up to $24$ and all
four token shapes, and this is what makes it the one forced-antipode context
that occurs.  The classification of the direct-edge two-input
row given above is of a different kind, since it is proved by hand and does
not depend on the search at all.

\section{Sequential realization}\label{sec:realization}

\begin{lemma}[realization]\label{lem:realization}
Let $M$ be the least common multiple of the catalogue denominators and of
$24$.  There is an absolute constant $A$ such that, for every decomposed core
with a family assigned to every owner as in Section~\ref{sec:cells} and
$n\ge 10^6(D+1)$, parameter values in $M\Z\cap[-A(D+1),A(D+1)]$ can be
chosen bottom-up, and the negatives of the tokens can be placed at stationary
vertices of the ordinary part of the tree, so that all special labels are
distinct and nonzero, $|S|\le 30D+3$, and the special support is closed under
negation up to at most one zero-sum triple whose negatives are unused.
\end{lemma}

Our proof chooses every parameter of every cell numerically, one coordinate
at a time, in post-order.  The tokens, the zero-sum triples that (E2) leaves
after the mirror pairs, are not cancelled against one another: the negatives
of each token are placed instead at a \emph{sink} in the ordinary part of the
tree, either directly in a zero-sum block of size three or at least five, or,
two tokens at a time, in a five-vertex \emph{absorber} built from stationary
vertices with two children (Lemma~\ref{lem:absorber}).  The hypothesis
$n\ge 10^6(D+1)$ is what guarantees enough sinks
(Lemma~\ref{lem:packing}), and a finite check over the catalogue shows that
the four extra labels of an absorber never coincide identically with a label
of either of its two source cells (Lemma~\ref{lem:compat}).  Thus no symbolic
parameter is ever live, apart from the forced-antipode children and the
two-input rows that have their own treatment, the number of forbidden
conditions at every choice is $O(|S|)$, and Lemma~\ref{lem:charge} gives boxes
of side $O(D+1)$.  The full argument is given in Appendix~\ref{app:realization},
together with the explicit value $A=\AConst$, which follows from the catalogue
constants (width at most $\NumMaxWidth$, at most $\NumParams$ parameters,
coefficients at most $\NumCoefBound$ over denominators dividing $12$) and
Lemma~\ref{lem:charge}; nothing is ever divided, so no lattice or coefficient
enlargement enters.  We have made no attempt to optimise it.  An earlier
version cancelled the tokens in pairs along the root path; why this cannot
give a linear threshold is recorded in Remark~\ref{rem:whysinks}.

\section{Dense completion and the proofs of Theorems~\ref{thm:main} and~\ref{thm:linear}}\label{sec:dense}

After Lemma~\ref{lem:realization} the unused labels form the set
$X=\Z_n\setminus(\{0\}\cup S)$, which is closed under negation up to at most
one leftover triple $T$ (the case of an odd number of tokens).  The negatives
$-T$ are placed in an odd ordinary block of size at least three, which exists
by parity whenever such a triple is left over and no block has size one; write
$r$ for that block's prescribed size and put
$X'=X\setminus(-T)$.  Since $r$ is odd and at least three, the remainder $r-3$
is zero or at least two, and the block is dropped from the list when it is
zero.  It remains to partition $X'$ into zero-sum blocks of the resulting
sizes, all at least two, with $|\Z_n\setminus X'|\le 2A(D+1)+4$.

\begin{lemma}[linear dense completion]\label{lem:dense}
There are absolute constants $\eta>0$ and $n_1$ such that for every odd
$n\ge n_1$, every symmetric $X\subseteq\Z_n\setminus\{0\}$ with
$|\Z_n\setminus X|\le\eta n$, and every list of block sizes $r_w\ge 2$ summing
to $|X|$, the set $X$ has a partition into zero-sum blocks of sizes $r_w$.
\end{lemma}

\begin{proof}[Proof of Lemma~\ref{lem:dense}]
Write $n=2K+1$.  Since $n$ is odd and $X$ is symmetric with $0\notin X$, the set
$X$ is a disjoint union of inverse pairs $\{j,-j\}$; let $J\subseteq[1,K]$ be
the set of magnitudes occurring, so that $|X|=2|J|$ and
$|\Z_n\setminus X|=1+2(K-|J|)$.  In particular $|X|$ is even, so the number
$o$ of odd block sizes is even.  Let $m_2$ be the number of blocks of size two
and put $N=\sum_{r_w\ge3}r_w$, so that $|X|=N+2m_2$.  We fix
$\eta\le1/20$ and argue by cases according to the size of $N$.

\emph{Case 1: $N\le n/100$.}  In this case we use no external input.  A
\emph{twin} is a pair of magnitudes $a<b$ in $J$ with $a+b\in J$; it yields the
two zero-sum triples $\{a,b,-(a+b)\}$ and $\{-a,-b,a+b\}$, which are disjoint
and whose union $\pm\{a,b,a+b\}$ is symmetric.  We choose $o/2$ twins with
pairwise disjoint magnitude sets, greedily.  At the start of step $s$ the
forbidden magnitudes are those outside $J$, of which there are at most
$\eta n/2$, together with the $3(s-1)\le 3o/2\le N/2\le n/200$ already used.
Hence at most $(\eta+1/100)K+2$ magnitudes are forbidden.  On the other hand
the number of triples $\{a,b,a+b\}\subseteq[1,K]$ with $a<b$ is
$\sum_{c=3}^{K}\lfloor(c-1)/2\rfloor\ge K^2/5$ for $K$ large, whereas a single
magnitude lies in at most $2K+K/2+3$ of them.  Since
$(\eta+1/100)\cdot(5/2)<1/5$ for $\eta\le1/20$, a permissible twin survives at
every step once $n$ is large, although the margin is not optimised.  Now pair
the odd blocks arbitrarily, give the two triples of the $s$-th twin to the two
blocks of the $s$-th pair, and distribute the remaining magnitudes as inverse
pairs, $p(r_w)=\lfloor r_w/2\rfloor-[r_w\text{ odd}]$ of them to the block
$w$.  The count is exact, because
$\sum_w p(r_w)=(|X|-3o)/2=|J|-3o/2$ is precisely the number of magnitudes left
over.  Every block then has its prescribed size and sum zero.

\emph{Case 2: $N>n/100$.}  Here we run the argument of
\cite[Lemma~6.27]{MuyesserPokrovskiy2025} in $\Z_n$, which has no involutions,
and with the missing pairs held out.  Put $p=N/(2|J|)$, so
$1/100\le p\le1$ because $2|J|=|X|\le n$.  Enumerate all $K$ inverse pairs of
$\Z_n$ and let $Y$ be a $p$-random set of indices; by
\cite[Lemma~6.26]{MuyesserPokrovskiy2025} applied to the partition of $\Z_n$
into these pairs together with $\{0\}$, the union of the selected pairs splits
as $Q\sqcup R\sqcup S$ with $Q,R$ disjoint $p/2$-random subsets of $\Z_n$ and
$S\subseteq\{0\}$.  Let $Y_J$ be the set of selected indices whose magnitude
lies in $J$; it is $p$-random on $|J|$ indices with integral mean
$p|J|=|J|-m_2$, so its median equals its mean and, with probability at least
one half, $|Y_J|\le|J|-m_2$, while Chernoff's bound gives
$|Y_J|\ge|J|-m_2-\sqrt{n}\log n$ with high probability.  Fix an outcome in
which both hold and in which $Q$ and $R$ satisfy the conclusions of
\cite[Lemmas~6.23 and~6.25]{MuyesserPokrovskiy2025}.  Reserve $m_2$ of the
$|J|-|Y_J|$ unselected available pairs for the blocks of size two, let $J^{*}$
be the union of the at most $\sqrt n\log n$ remaining unselected available
pairs, and put $Z=J^{*}\cup\bigcup_{i\in Y_J}\{i,-i\}$, a union of inverse
pairs with $|Z|=2|J|-2m_2=N$ and $\sum Z=0$.

Reduce the multiset $\{r_w:r_w\ge3\}$ to a multiset $M'\subseteq\{3,4,5\}$ as
in the first paragraph of the proof of
\cite[Lemma~6.25]{MuyesserPokrovskiy2025}, recombining at the end, and split
$M'=M^{Q}\sqcup M^{R}$ by giving $M^{Q}$ the floor and $M^{R}$ the ceiling of
half the multiplicity of each of $3,4,5$.  The multiplicities need not be even,
so the two parts need not have equal sums; however they differ in multiplicity
by at most one in each of the three sizes, hence
$\bigl|\sum M^{Q}-\sum M^{R}\bigr|\le12$, and a bounded discrepancy is
absorbed by the slack reserved below.

Write $\delta=(p/2)^{10^{10}}/\log(n)^{10^{24}}$ for the perturbation radius
that \cite[Lemma~6.25]{MuyesserPokrovskiy2025} allows at density $p/2$, and
let $\varepsilon=\delta/24$, which lies in the range permitted by
\cite[Lemma~6.23]{MuyesserPokrovskiy2025} once $n$ is large.  Both
$Q\setminus Z$ and $R\setminus Z$ are contained in $\{0\}\cup(\Z_n\setminus X)$
and hence have size at most $|\Z_n\setminus X|$, and
$\bigl|\sum M^{Q}-pn/2\bigr|\le|\Z_n\setminus X|+O(\sqrt n\log n)$.  Assume
$|\Z_n\setminus X|\le\delta n/48$, which is the hypothesis of the lemma with
$\eta=\delta/48$.  The two hypotheses of
\cite[Lemma~6.23]{MuyesserPokrovskiy2025} are then met and not merely its
conclusion is used: $|Q\setminus Z|\le|\Z_n\setminus X|\le\varepsilon n/2$,
and $\bigl|\sum M^{Q}-pn/2\bigr|\le\varepsilon n/2+O(\sqrt n\log n)
\le\varepsilon n$ once $n$ is large.  Then
\cite[Lemma~6.23]{MuyesserPokrovskiy2025}, applied with
$R=Q$, $Z$ as above, $m=\sum M^{Q}$ and $g=0$, returns $Q'\subseteq Z$ with
$|Q'|=\sum M^{Q}$, $\sum Q'=0$ and
$|Q'\triangle Q|\le6\varepsilon n=\delta n/4$.  Put $R'=Z\setminus Q'$;
then $|R'|=\sum M^{R}$, and $\sum R'=\sum Z-\sum Q'=0$, and
$$|R'\triangle R|\le6\varepsilon n+|\Z_n\setminus X|+|J^{*}|
  \le\tfrac{\delta n}{4}+\tfrac{\delta n}{48}+O(\sqrt n\log n)\le\delta n$$
once $n$ is large, and likewise $|Q'\triangle Q|\le\delta n/4\le\delta n$.
Thus both perturbations stay strictly inside the allowance, which is the point
of taking $\varepsilon$ a twenty-fourth of $\delta$ rather than a sixth, while
the deletion radius $\eta=\delta/48$ is what keeps the two \emph{input}
hypotheses inside $\varepsilon n$; the two budgets are separate and an earlier
version of this proof conflated them.
Therefore \cite[Lemma~6.25]{MuyesserPokrovskiy2025} applies to $Q'$ with the
reference set $Q$ and the multiset $M^{Q}$, and to $R'$ with $R$ and $M^{R}$,
and returns zero-sum partitions of $Q'$ and of $R'$ of the prescribed sizes.  Together with
the $m_2$ reserved inverse pairs these fill every block, since
$Z\sqcup(\text{reserved pairs})=X$.  When $m_2=0$ one has
$p=1$, $Y_J$ is all of $J$ and $J^{*}$ is empty, and the argument degenerates
to a direct application of \cite[Lemma~6.25]{MuyesserPokrovskiy2025} to $X$.
Thus $X$ has the required partition
whenever $|\Z_n\setminus X|\le\delta n/48$, and $\delta$ depends only
on $p\ge1/100$ and on $n$ through the logarithm.
\end{proof}

\begin{remark}[the two deletion bounds of \cite{MuyesserPokrovskiy2025}]\label{rem:mplinear}
Two readings of \cite{MuyesserPokrovskiy2025} enter through the parameter
$\varepsilon$ of the last proof, and they give the two thresholds of
Section~\ref{sec:intro}.  The printed perturbation radius of their
Theorems~4.3--4.7, and hence of their Lemmas~6.23 and~6.25, is
$p^{A}n/(\log n)^{A}$; since the proof above uses those lemmas only with
$p\ge1/200$, this gives Lemma~\ref{lem:dense} with $\eta n$ replaced by
$n/(\log n)^{A}$, and hence Theorem~\ref{thm:main}.  Their Section~7, under
\emph{Bounds in the main theorem}, records that $\log n$ may be replaced by
$\log|G'|$ throughout the proof, the only source of logarithms being their
short commutator-product theorem, and for an abelian group the commutator
subgroup is trivial and the commutator product is empty; the authors record
the consequence $g(n)=\Omega(n)$ for cyclic groups in their Problem~7.4.  One
should read their commutator length as bounded by $\max(1,\log_4|G'|)$ rather
than divide by $\log1$.  Their Lemmas~6.23 and~6.25 are deduced from
Theorems~4.4--4.6 with no further logarithmic loss, the auxiliary error terms
being of order $\sqrt n\log n$, so at every fixed reservoir density this
reading leaves a positive constant perturbation radius and gives
Lemma~\ref{lem:dense} as stated, which is what Theorem~\ref{thm:linear} uses.
The authors describe this improvement as requiring essentially no
modification of their proof and state its consequence for cyclic groups,
$g(n)=\Omega(n)$, in their Problem~7.4; neither is a numbered theorem in
their paper, and we have asked the authors to confirm the linear form of
Lemmas~6.23 and~6.25.  The constant $\eta$ is in either case inherited from
their argument and is not effective.
\end{remark}

\begin{proof}[Proof of Theorems~\ref{thm:main} and~\ref{thm:linear}]
Root $T$ at a branch vertex, decompose by Section~\ref{sec:decomp}, assign
families by Section~\ref{sec:cells}, realize by Lemma~\ref{lem:realization}
with $|S|\le 30D+3$ and magnitudes at most $A(D+1)$ (the hypothesis
$n\ge 10^6(D+1)$ of that lemma is implied by $n>C(D+1)$ once
$C\ge 10^6$), and place the leftover triple if there is one.  With the linear
deletion bound of Remark~\ref{rem:mplinear}, Lemma~\ref{lem:dense}
completes the labelling as soon as $2A(D+1)+1\le\eta n$; every stationary
vertex then has a zero-sum block, every special vertex has the prescribed sum,
and Lemma~\ref{lem:sigma} gives the bijection, so $C\ge 3A/\eta$ and
$n_0\ge n_1$ prove Theorem~\ref{thm:linear}.  With the printed perturbation
radius of \cite{MuyesserPokrovskiy2025} in place of the linear one the same
argument applies when $2A(D+1)+1\le n/(\log n)^{A'}$, which proves
Theorem~\ref{thm:main}; its final assertion follows by taking $D\le k$ and
$n$ large.
\end{proof}

% Section draft: the elementary route to a linear threshold.
% To be \input after the dense-completion section once the constant is settled.

\section{An elementary route to the linear threshold}\label{sec:linearroute}

Theorem~\ref{thm:main} rests on the abelian reading of
\cite{MuyesserPokrovskiy2025} through Lemma~\ref{lem:dense}.  In this section
we record how far our reduction removes that dependence by elementary means.  The
outcome is that Lemma~\ref{lem:dense} reduces to a single statement about
permutations of an interval, which we prove in the two extreme cases and in a
family of composite cases, and which we have verified exactly for every order
up to $30$.  We state it as a conjecture; a proof would make
Theorem~\ref{thm:main} independent of \cite{MuyesserPokrovskiy2025} and would
give an explicit value of $C$.

\subsection{From the residual set to Heffter's difference problem}

After Lemma~\ref{lem:realization} the special support is $\pm I$ with
$I\subseteq[1,M]$ and $M=A(D+1)$, and the residual set is
$X=\pm([1,K]\setminus I)$ in $\Z_n$, $n=2K+1$.  In the branch of
Lemma~\ref{lem:dense} that needs the completion theorem, $X$ has to be
partitioned into signed zero-sum triples $\{\pm a,\pm b,\pm c\}$.  Writing a
triple by its three magnitudes $a<b<c$ in $[1,K]$, the two possibilities are
$a+b=c$ and $a+b+c=n$; the first is a pair $\{b,c\}$ at distance $a$ and the
second is the pair $\{b,-c\}$ at cyclic distance $a$ in $\Z_n$.  For
$I=\emptyset$ this is Heffter's first difference problem when $n\equiv1\pmod
6$, solved for every $n$ by Peltesohn~\cite{Peltesohn}, and the general case
asks for the same partition of $[1,K]\setminus I$, that is, for Heffter's
problem with the small differences $I$ forbidden.  The hardest case in our experiments is
$I=[1,M]$, and it is the one we analyse.

\subsection{Folded sum permutations}

Let $K-M=3t$.  In every solution our searches found, the $t$ smallest elements
$M+1,\dots,M+t$ can be taken as the distances $a$, and the remaining $2t$
elements split into a lower half $M+t+1,\dots,M+2t$ and an upper half
$M+2t+1,\dots,K$ with each pair $\{b,c\}$ taking $b$ from the lower and $c$
from the upper half.  Writing $a=M+1+i$, $b=M+t+1+j$, $c=M+2t+1+k$ with
$i,j,k\in[0,t-1]$ and $s=M+1-t$, the two triple types become
\[
  c-b=a\iff k=i+j+s,\qquad b+c=n-a\iff k=(2t-1)-(i+j+s).
\]
So a solution is a permutation $\pi$ of $[0,t-1]$, $j=\pi(i)$, such that the
numbers $u_i=i+\pi(i)+s$ lie in $[0,2t-1]$ and their folds
$\min(u_i,\,2t-1-u_i)$ form a permutation of $[0,t-1]$.  We call such a $\pi$
a \emph{folded sum permutation} of order $t$ and shift $s$, and write
$\mathrm{FSP}(t,s)$.  Equivalently, with $\tau_i=i+\pi(i)$, the $\tau_i$ are
distinct, lie in the window $W=[-s,2t-1-s]$, and no two of them sum to the
reflection constant $c=2t-1-2s$ of $W$; that is, the sum set of $\pi$ is a
transversal of the reflection of $W$.  Since $\sum_i\tau_i=t(t-1)$ for every
permutation, the mean of the transversal exceeds the centre of $W$ by
$s-\tfrac12$.

The range $3M+1\le K\le 7M+3$ of Lemma~\ref{lem:dense} corresponds to
$(1-t)/2\le s\le(t+1)/2$, and beyond $K\ge7M+3$ Simpson's theorem on Langford
sequences applies with no folding at all.

\begin{conjecture}\label{conj:fsp}
For every $t\ge5$ and every integer $s$ with $(1-t)/2\le s\le(t+1)/2$ there is
a folded sum permutation of order $t$ and shift $s$.
\end{conjecture}

The window is sharp: outside it no $\mathrm{FSP}(t,s)$ exists, because the
$s$ largest values of $W$ (for $s\ge1$) have no reflection partner in the
range of the $\tau_i$ and must all be attained, which is impossible once
$s>(t+1)/2$.  Our constraint programs verify the conjecture exactly, for every $3\le t\le30$ except $t=4$, where $s=-1$ and $s=2$
fail; unlike Simpson's theorem there are no congruence exceptions inside the
window, the folding removes them.

\subsection{What is proved}

\begin{proposition}\label{prop:fspcases}
Conjecture~\ref{conj:fsp} holds in the following cases.
\begin{enumerate}
\item $s\in\{0,1\}$: the identity is an $\mathrm{FSP}(t,s)$.
\item $t$ odd and $s\in\{(1-t)/2,(t+1)/2\}$: the rotation
  $\pi(i)=i+(t-1)/2 \bmod t$ is an $\mathrm{FSP}(t,s)$ at both endpoints.
\item (Composition.)  If $a$ is odd and $\sigma$ is an
  $\mathrm{FSP}(b,s')$, then
  \[
    \pi(au+w)=a\,\sigma(u)+\bigl((w+\tfrac{a-1}{2})\bmod a\bigr),\qquad
    u\in[0,b-1],\ w\in[0,a-1],
  \]
  is an $\mathrm{FSP}(ab,s)$ with $s=a(s'-1)+(a+1)/2$.
\end{enumerate}
\end{proposition}

\begin{proof}
(1) The sums $2i$ are distinct and even, the constant $c$ is odd, and
$0\le 2i\le 2t-2\le 2t-1-s$ exactly when $s\le1$; the lower bound $2i\ge-s$
needs $s\ge0$.

(2) For the rotation the sums $i+\pi(i)$ are $2i+(t-1)/2$ for
$i<(t+1)/2$ and $2i-(t+1)/2$ for $i\ge(t+1)/2$, which together form the
interval $[(t-1)/2,(3t-3)/2]$.  At $s=(1-t)/2$ this interval is the lower half
of $W$, at $s=(t+1)/2$ it is the upper half, and a half of $W$ contains no
reflected pair.

(3) On each block the inner rotation has sums $w+\rho(w)$ forming the interval
$I=[(a-1)/2,(3a-3)/2]$, which is carried to itself by $x\mapsto 2a-2-x$.  The
sums of $\pi$ are $a(u+\sigma(u))+I$.  Put $s=(a+1)/2+aq$.  The condition
$a(u+\sigma(u))+I\subseteq[-s,2ab-1-s]$ is $u+\sigma(u)\in[-(q+1),2b-q-2]$,
the window of $\mathrm{FSP}(b,q+1)$; and the reflection constant
$c=2ab-1-2s=a(2b-2q-1)-2$ sends $a\,y+x$ to $a(2b-2q-3-y)+(2a-2-x)$, so it
reflects $I$ onto itself and the block index $y$ by the constant
$2b-1-2(q+1)$ of $\mathrm{FSP}(b,q+1)$.  Therefore the sum set of $\pi$ is a
transversal exactly when the sum set of $\sigma$ is, and $s'=q+1$.
\end{proof}

Case (2) has two readings: at the lower endpoint it is a hooked Langford
sequence of defect $(t+1)/2$, at the upper endpoint the columns
$(i,\pi(i),c-i-\pi(i))$ form a $3\times t$ magic rectangle.  Case (3) covers
the residues $s\equiv(a+1)/2\pmod a$ for every odd divisor $a$ of $t$; it does
not reach prime $t$, and we have not found an additive counterpart.

\subsection{Consequences}

If Conjecture~\ref{conj:fsp} holds then the clean tail $[M+1,K]$ has a
maximum packing into difference triples for every $K\ge3M+1$ (the cases
$K-M\not\equiv0\pmod3$ reduce to it by omitting the one or two smallest
elements, and the finitely many cases $K\le3M+3$ are magic rectangles with one
or two empty cells, verified for $M\le30$).  Together with the proved branch of
Lemma~\ref{lem:dense} for $K\ge7M+3$, this gives Lemma~\ref{lem:dense} with
$\eta=1/(14A+O(1))$, hence Theorem~\ref{thm:main} with an explicit constant
$C$ and without any use of \cite{MuyesserPokrovskiy2025}.  The punctured
cases $I\subsetneq[1,M]$ are the same object with the extra small elements
as additional distances, and were found feasible for every random $I$ we
tried; we have not reduced them to Conjecture~\ref{conj:fsp} in general.


\section{Computations and certificates}\label{sec:verification}

Three programs enter the proof, and we describe them in the order in which
their results are used.

\paragraph{The symbolic search.}
For a context the search builds the label forms of Section~\ref{sec:cells}
with integer coefficients bounded by $24$ over a common denominator $L\le 12$
(the families built by the two identities of Section~\ref{sec:decomp} are not
subject to this cap, and the largest coefficient anywhere in the catalogue,
$\NumCoefBound$, is attained by the equal-pair blocks; it is this
catalogue-wide value that enters Lemma~\ref{lem:realization}),
encodes the $\PO$ identity as a permutation of the output multiset onto the
label multiset (one Boolean per pair, each true Boolean equating two
coefficient vectors), encodes the signature (E2), the distinctness (E3), the
token rank condition (E4) and the head freedom (E5) as further Boolean and
linear constraints, and asks CP-SAT~\cite{ORTools} for a feasible assignment.
The same model with the cleanliness conditions (E6) is tried first, and the
alternatives (token of rank three, token of rank two, no dedicated token
parameters, two tokens) are tried in that order.  Every feasible assignment
is stored, and the menu of a context lists all of them in order of
preference, families containing the antipode of their own head last.  For
the contexts with inputs the model is run a second time with the shared edge
required to be a token member and the antipodes of the inputs forbidden
(and a third time, for two-input cells, with both inputs required to be
token members).  The resulting \emph{variant} families are kept in separate
tables and are the ones used when a child is forced-antipode.  Each context is tried under nine settings of (denominator, number of
internal parameters, coefficient bound), namely $(1,1,3)$, $(2,1,3)$,
$(3,1,3)$, $(1,2,4)$, $(2,2,4)$, $(3,2,4)$, $(4,2,4)$, $(6,2,6)$ and
$(12,2,12)$, with the head multiplier free, and with a time limit between
$150$ and $900$ seconds per context and setting on a $128$-core workstation.
In all, \NumAttempted{} contexts with widths up to $29$ labels were attempted,
in \NumSolverRuns{} runs of the solver, and \NumReceived{} of them received at
least one family.
The search is exact in the sense that an infeasible model is a proof of
non-existence within the coefficient and denominator bounds, and it is in
this sense that the impossibilities of Section~\ref{sec:cells} are stated.

\paragraph{The independent checker.}
The checker reads only the JSON record of a family, rebuilds the cell
topology from the row specification (head chain, arms, ports, active and
input children, root row), recomputes every output as an exact rational
combination of the parameters, and verifies (E1) by matching the two
multisets of coefficient vectors, (E2)--(E5) directly, and the numerical
identity at three random integer points, and it shares no code with the search.
All \NumCatFamilies{} families of the catalogue and all
\NumVariantFamilies{} variant families pass, as do the families obtained
from them by one six-edge insertion on every arm and by one twelve-label
insertion on every head chain.  The command
\begin{verbatim}
python3 scripts/check_alphabet_families.py data/alphabet_families.json
\end{verbatim}
re-verifies the whole catalogue in about two minutes on a laptop and prints
one line per context.

\paragraph{The gate and the closure check.}
Two adjacent cells can fail to be compatible in only one way.  A label of the
parent that depends on its input alone is $qx$ for a rational $q$, a label of
the child that depends on its own head alone is $q'y$, and if the two ratios
agree once the shared edge is identified then the two labels are equal as
forms and no choice of parameters separates them.  Our gate enumerates every
ordered pair of contexts, tries every pair of families from the two menus, and
reports the context pairs for which no combination avoids such a coincidence.
Since the sink route places the negative of every unit label elsewhere in the
tree, the gate also checks the negative of each unit label of either cell
against the labels of the other (this is what Lemma~\ref{lem:compat} needs
for the cells that violate the second clause of (E6)); the addition changes
the count by one pair, which is not realizable.
Over the catalogue it reports \NumGateFail{} ordered pairs.  However, most
ordered pairs of contexts cannot occur in a tree at all, and restricted to the
\NumPairs{} ordered pairs that the enumeration of Lemma~\ref{lem:contexts}
records it reports \NumGateFailRealizable{}.  All but one of these have the
forced-antipode child of Section~\ref{sec:cells} as their child and are
therefore absorbed by the rule of that section and not by the gate.  The
exception appeared only when the enumeration was run at three active children,
and it is the pair whose parent is the root context with a lone arm of
residue two and one active child, and whose child is the bottom cell with arms
of lengths one and two.  Here the parent's forced ratios are $-1$, $-\tfrac12$
and $\tfrac12$ in every family of its menu, and the child carries the antipode
of its own head in every family of its menu, so the two collide at every head
multiplier the search produces.  However, the gate is a sufficient condition
and not a necessary one: we join the child into the parent and solve the two
label sets together as a single cell of width ten, for which the catalogue
carries a family, and an identity between labels of one cell cannot arise.  We
note that the joining stops there, since the parent is the root and the child
is a bottom cell, so the enlarged cell has neither a parent of its own nor an
active child, and this is the same device as the joint of the absorption
rule.  A second
program checks the conditions that the absorption rule needs, namely that
every context that occurs has a family, that every context with a free active
child has an input-in-token family or a joint that is not forced-antipode,
that the folded letter has such a joint, and that every two-input context
either has a family with both inputs in tokens or has a port joint carrying
the second pinned child in a token.  We note that our scheduler deliberately
stops one step short of the absorption rule, since it emits the
forced-antipode cell as an owner of its own and does not join it into its
parent, because which of the two routes applies depends on the certificate
table and not on the tree.  The context list therefore carries the joint of
every free-child context and the port joint of every two-input context, and
the second program is what verifies that the rule closes over them.  These
conditions are a finite list and the program reports them one by one, so that
the state of the verification is visible at any time.  All of them hold: every
one of the \NumEmittable{} contexts that occur has a family, all
\NumFreeChildContexts{} free-child contexts absorb a forced-antipode child,
the folded letter and all $40$ two-input contexts do the same, and the gate
leaves nothing outside the absorption rule and the single join above.

\paragraph{The scheduler and the constructor.}
Our scheduler implements the rules of Appendix~\ref{app:contexts} and returns,
for a rooted tree, the list of owners with their contexts.  It is the
regression for Lemma~\ref{lem:contexts} and was run on all labelled trees with
at most nine vertices and on $169{,}000$ random trees of odd orders between
$11$ and $401$.  It roots the core at a branch vertex, as
Lemma~\ref{lem:contexts} requires; an earlier version of it took the root from
its caller without checking the degree, which is what produced the only
uncovered owners we have seen.  Our constructor
implements Lemma~\ref{lem:realization} on top of the scheduler and the
catalogue, produces an integer labelling, reduces it modulo $n$, and checks
the bijection of vertex sums directly.  Each run writes a certificate listing
the owners, the chosen families, the parameter values and the final
labelling, so that a labelling can be re-verified without the constructor.
The constructor has two modes.  In its default mode it follows the
decomposition but completes the ordinary blocks by a direct search, and in
the sparse regime ($n\in\{51,101,201\}$ with $D\approx n/25$) it succeeded
on $116$ of $120$ random trees.  Its \emph{proof-faithful} mode follows the
earlier version of Lemma~\ref{lem:realization}, in which tokens were cancelled
in pairs along the root path (numeric coordinates, live carriers, deferred
pairing, the reservation of forced companions), and reports every departure
from that lemma as a named gap.  On the same $120$ trees it succeeded on $98$.
Of the twenty-two remaining runs, fourteen report a modular collision between
two special labels, or a pending triple whose negatives are already used
modulo $n$, seven report a failure of the completion step, and one reports
that no odd block was available for the singleton bridge.  On a comb attached
to a dense part, however, the same mode records the gap ``a third token became
pending'' at every second core vertex of the comb, with head labels doubling
along it, which is the failure of the pairing route recorded in
Remark~\ref{rem:whysinks}.  A third mode (\file{sink\_route/sink\_constructor.py})
implements Appendix~\ref{app:realization} as it now stands: it reserves the
sinks before any label is chosen, chooses every parameter numerically, fills
the sinks and completes the ordinary blocks.  It labels the comb attached to a
root with $12k$ stationary children for $k\in\{5,9,15\}$, the comb attached to
a complete binary tree of depth $7$, $8$ and $9$ for the same three values of
$k$ (using $39$ absorbers and $9$ direct sinks in all), and fifteen random
sparse trees of orders $101$, $201$ and $401$, with no lemma gap on any of
these $27$ trees; the largest label magnitude is below $7(D+1)$ and
$|S|<4.5(D+1)$, against the bounds $A(D+1)$ and $30D+3$ that we prove (with
units in place of tokens the same trees give the same outcomes, four of the
sparse trees now sending a quad to a block of size four, and the quad gadget,
the sink for the second token of a forced-antipode family and the exclusion
of a parent family above a residue-two ray were each exercised on a small tree
built for the purpose).  Whether a tree has a cancelling triple at all depends
on the root: none of the trees above has one at the root the constructor
chooses first, so the triples are exercised separately.  On the sparse tree of
order $201$ rooted at the vertex which does produce one, the three heads
$360$, $-4320$, $3960$ sum to zero, their negatives are placed together in one
ordinary block, and the labelling verifies
(\file{sink\_route/regression\_cancel\_triple/}).  On the
$293$ free trees of orders $7$, $9$ and $11$ it succeeds on $243$; of the fifty
$293$ free trees of orders $7$, $9$ and $11$ it succeeds on $243$; of the fifty
failures, $33$ are modular collisions, nine report that the tree has no sink at
all (no block of size three or at least five and no path of five stationary
vertices with two children), which is the hypothesis of
Lemma~\ref{lem:realization} failing, as it must at these orders, and eight fail
in a cell or in the completion.  Every successful certificate was re-verified
independently of the constructor.  The same mode realizes the two-level joints
of Lemma~\ref{lem:reset} (\file{sink\_route/reset\_joint\_constructor.py}): on
the comb of Appendix~\ref{app:realization} whose every core vertex has two
active children and one arm of length one, attached to a bushy root with
$n\approx 12.5(D+1)$, chains of $4$ to $8$ rows are cut into $2$, $2$, $3$,
$3$ and $4$ joints and every tree is labelled with the largest magnitude below
$6(D+1)$, whereas without the joints the same trees either fail or reach
$21(D+1)$; the scheduler's owner lists on all trees of order at most nine are
unchanged by the addition.  We note that all these runs use the dense
completion only through a direct construction of the leftover blocks (from the
antipodal pairs, which the sink route makes available, with the CP-SAT search
as a fallback), since at these orders the asymptotic Lemma~\ref{lem:dense}
does not apply, and their purpose is to confirm the bookkeeping of the cells
and of the sinks and not the theorem itself.  A modular collision between two
special labels is exactly the event that the hypothesis $n>C(D+1)$ excludes,
and at $n=201$ with $D\approx 8$ the hypothesis is very far from holding.

\paragraph{The absorber gate and the sinks.}
Lemma~\ref{lem:compat} needs, for every pair of cells that can share a
parameter, a role order of the two tokens for which none of the four absorber
forms coincides identically with a label of either cell, with the negative of
one, with another absorber form or with zero.  Our program
\file{sink\_route/absorber\_gate.py} checks this over the joint parameter space
of the two families, the child's head identified with the parent's input, for
every ordered pair of a catalogue family with one input and a token ($625$
families) and a catalogue family with a head and a token ($2{,}113$ families),
one family per context, namely the one the scheduler prefers there:
of the $1{,}320{,}625$ pairs none has fewer than eight collision-free orders
among the thirty-six, and $1{,}077{,}257$ have all thirty-six.  The same check
over the $4{,}242{,}904$ pairs with a two-input parent (the child attached at
either input, either token of the parent), the $574{,}907$ pairs in which at
least one side uses its second token, and the $401$ families carrying two
tokens (both tokens of one cell as the two sources) finds no pair without a
collision-free order either, the minima being six, eight and three orders
respectively (\file{sink\_route/absorber\_gate\_ext.py}).  The same test was
then run over \emph{every} entry of every menu rather than one family per
context (\file{sink\_route/full\_menu\_absorber\_gate.py}): $2{,}175$ parent
families against $4{,}431$ child families, $9{,}637{,}425$ ordered pairs, again
no pair without a collision-free order, the minimum being six.  A second
implementation, which carries the labels as exact rationals over an explicit
joint coordinate list and shares no code with the first, reproduces that run
pair for pair, including the whole histogram of collision-free orders
(\file{sink\_route/verify\_full\_menu\_gate.py}); the certificate records the
catalogue's SHA-256 digest.  Thus the pairing of
Lemma~\ref{lem:compat} may be arbitrary.  The units themselves were scanned
over the $6{,}975$ token-bearing families of all menus: $528$ contain an
antipodal pair among their token labels, of which $394$ leave a quad and
$134$ leave nothing to place, and every family with $-x$ in a token also has
$x$ in a token, so the case the units were introduced for never arises in
any other form.
For the residue-two ray, exactly one single-input family
(\texttt{h4\_6-6\_p0\_act}, first of its menu) has a label $\pm x/3$ or
$\pm 2x/3$ and its context has another family without, and one two-input
joint family (\texttt{h3\_1-3-3-3\_p0\_A2}) has such a label on one of its
two inputs, to which the constructor never assigns a ray; no context is left
without an admissible family.  The counting of Lemma~\ref{lem:packing} was also
run on real schedules (\file{sink\_route/packing\_statistics.py}): on random
sparse trees with $D\approx n/25$ the direct capacity alone exceeds the number
of owners by a factor of at least twelve, on a comb of $k$ core vertices
attached to a root with $12k$ stationary children the root's block absorbs
everything, and on a comb attached to a complete binary tree of depth $d$,
where no block has size three or more, the greedy packing finds $2^{d-4}$
disjoint absorbers, enough for every $k\le 15$ from $d=7$ on.  These runs are
evidence for the constants, not part of the proof.

\section{Discussion}\label{sec:discussion}

Previous work on edge-graceful trees is of two kinds.  Explicit labellings
are known for several families, for example paths, stars, spiders and brooms
(see \cite{GallianDS6} for the list; the regular spiders of
\cite{CabanissLowMitchem} already carry unboundedly many vertices of degree
two, although in a single prescribed shape), and the zero-sum partition theorem of
\cite{KaplanLevRoditty2009} covers every odd-order tree with at most one
vertex of degree two.  Our earlier work \cite{Meng2026two} extended that
theorem to two such vertices with explicit (and effective) constructions.
Our result lies between the two: it covers all trees
in which the vertices of degree two are a small fraction of the order, with
no restriction on where they sit, at the price of an ineffective constant.

The constant $C$ is not effective because Lemma~\ref{lem:dense} is not,
whereas the constant $A$ of Lemma~\ref{lem:realization} is explicit
($A=\AConst$).  An effective version of the zero-sum matching with
linear deficiency would therefore make Theorem~\ref{thm:linear} effective
without any change to the construction, and Section~\ref{sec:linearroute}
reduces this to Conjecture~\ref{conj:fsp}.  However, the degree-two vertices
remain a bottleneck in a stronger sense: the special support has size linear
in $D$ because every degree-two vertex is non-stationary, so our argument
cannot reach trees in which these vertices have positive density.  In
particular it says nothing about caterpillars with many subdivided edges or
about spiders with long legs, for which the explicit constructions in the
literature remain the only tool.  Whether the density threshold can be
raised, e.g.\ to $D=o(n)$ or to $D\le cn$ for a small absolute $c$, seems to us
the natural next question.  The local families are indifferent to the density,
so the obstacle is entirely in the completion step, where the leftover set
would no longer be nearly the whole group.

Two further remarks concern the computer-assisted part.  First, our
catalogue is finite by design and every family in it is a rational identity
in a few parameters, and an independent re-verification is for this reason
possible at all.  The search is only a way of finding the identities, and the
proof does not depend on the search having been exhaustive except in the four
impossibility statements of Section~\ref{sec:cells}, each of them a
finite enumeration, and one of the four structural facts we use, the
classification of the direct-edge two-input row, is proved by hand.  Second, although our scheduler's regression covers all
small trees, it is evidence for, not a proof of, Lemma~\ref{lem:contexts}.
The proof is the case analysis of Appendix~\ref{app:contexts}, and the
scheduler is our way of keeping that analysis honest.

\section{AI-assisted research statement}\label{sec:ai}

Anthropic's Claude, OpenAI's ChatGPT, Google's Gemini and Moonshot AI's Kimi
were used as coding and analysis assistants throughout this project.  They
implemented and executed the processing, reasoning, and statistical code,
prepared documentation, surveyed literature, and assisted in preparing
manuscript text.  The author assumes responsibility for all content.

\section{Data and code availability}\label{sec:data}

The search (\file{symbolic\_free\_token\_cpsat.py}), the batch driver, the
independent checker (\file{check\_alphabet\_families.py}), the scheduler
(\file{context\_scheduler.py}), the abstract enumeration
(\file{enumerate\_emittable\_contexts.py}), the gate
(\file{check\_lookahead\_gate.py}), the closure check
(\file{closure\_check.py}), the literal test of (E4) and (E6)
(\file{audit\_e4\_e6.py}, \file{e6\_coverage.py}), the non-constancy check of
the absorber gate and the sink statistics (\file{sink\_route/absorber\_gate.py},
\file{sink\_route/full\_menu\_absorber\_gate.py} with its
independent verifier,
\file{sink\_route/packing\_statistics.py}), the constructors (including
\file{sink\_route/sink\_constructor.py} with its regression driver), the context
list, the certificate catalogue (\file{alphabet\_families.json}) with its
variant tables, and the certificates of the runs reported in
Section~\ref{sec:verification} are available at
\url{https://github.com/lsmeng/sparse-edge-graceful} and archived at
\url{https://doi.org/10.5281/zenodo.22287387} (concept DOI, resolving to the
latest version).

\appendix
\section{Context completeness}\label{app:contexts}

We prove Lemma~\ref{lem:contexts} by describing the assignment step by step.
The executable scheduler follows the same steps and is our regression for the
case analysis.  Process the core bottom-up and let $v$ be a nonroot core
vertex with parent core edge of length $h$, arm multiset $A$, $k$ active
children (children which are owners) and $q$ ports.  Since $v$ has degree at
least three in $T$ we have $|A|+k+q\ge 2$.

\paragraph{Step 1 (arms).}
Lemma~\ref{lem:units} partitions $A$ into neutral units and a residual $R$ of
one of the six types.  The units are $\PO$ with zero head sum and with support
closed under negation, so they do not enter the cell of $v$.  However, when a
unit is needed at $v$ itself (Steps~3 and~4) one of them is retained and its
arms are added to the residual.

\paragraph{Step 2 (active children).}
Cancelling a set of $j\ge 2$ active children means prescribing their heads to
sum to zero: for $j$ even we use $j/2$ mirror pairs, and for $j$ odd one free
zero-sum triple together with $(j-3)/2$ pairs.  These heads are free parameters
of the children's families by (E5), with one exception, and a child of the
joined types, whose head is forced, is paired with a sibling of free head, the
sibling taking the negative of the forced head.  The exception is a pair of two
children that both have forced heads, which our scheduler does produce.  It is
also arrangeable, one level further down: a child whose head is forced in this
way has context $(0,\{1\},0,\mathrm{free})$, its own active child is of free
type, and the head of that grandchild is a free parameter by (E5), so setting
the grandchild's head makes the two forced heads negatives of one another.  We
note that this uses (E5) at the grandchild and not at the child, and that it
does not recurse further, since the grandchild is of free type and is therefore
not one of the rigid rows.

Whether all of the active children can be cancelled depends on what Steps~3
and~4 can then do with the residual.  If cancelling all of them leaves $v$
covered, we cancel all of them, and $v$ has at most one active child left,
whose head is the input $x$.  Otherwise $v$ must spend active children on the
residual, and the rule is by the parity of their number $k$.  For $k$ even we
keep two and cancel $k-2$, so that $v$ is a \emph{two-input} owner with inputs
$x_1,x_2$; for $k$ odd we keep one and cancel $k-1$.  In both cases the number
cancelled is at least two whenever it is positive, so the cancellation above
applies.  At the root with a lone residual arm there is one deviation: $k=2$
keeps both, because keeping one would leave a single child to cancel, and every
$k\ge 3$ keeps one.

The configurations in which the residual cannot be covered without active
children, so that this rule is used, are the following, and they are the exact
condition our scheduler tests.  Off the root, $v$ has no ports, its residual is
a single arm, and no neutral unit was retained; then the residues modulo six
that need an active child are $1,2,4,6$ when the head chain has length $0$
modulo six and all of $1,\dots,6$ otherwise.  At the root with a lone residual
arm of residue $r$ and $q$ ports, the recruited port count of Step~5 has no
feasible value exactly when $q=0$ and $r\in\{2,4\}$, or $q=1$ and
$r\in\{1,3,5\}$; for $q\ge2$ there is always one.  At the root there is one
further case, the opposite-parity residual $\{1,2\}$ with no ports.  We note
that an earlier version of this step named only the portless residuals $\{1\}$
and $\{2\}$ and only $k=2$; that was too narrow, as a tree of order $19$ on
which our scheduler keeps two active children at a residual of residue four
shows.  With a subdivided parent
edge such cells have ordinary families, e.g.\
$(y,-y-x_1-x_2;\,-x_1-x_2,\,y+x_1+x_2;\,x_1,x_2)$ for $h=1$ and $R=\{1\}$ with
token $\{x_1,x_2,-x_1-x_2\}$, some of them with two tokens.  With a direct
parent edge the exact search shows that no family with a free head exists,
and $v$ uses instead one of the two \emph{rigid} rows
$(-x_1-x_2;\,-x_2,\,-x_1;\,x_1,x_2)$ for $R=\{1\}$ and
$\bigl(-\tfrac32(x_1+x_2);\,-\tfrac12(x_1+x_2),\,-(x_1+x_2),\,\tfrac12(x_1+x_2);\,x_1,x_2\bigr)$
for $R=\{2\}$, whose head is forced by the inputs and which are treated like
the residue-two ray in Appendix~\ref{app:realization}.  In our random-tree
regression this pattern occurs at roughly one to two percent of the vertices.

\paragraph{Step 3 (empty residual).}
Suppose $R=\emptyset$.  If no active child remains, $v$ is stationary and its
ordinary block consists of its $q$ ports, which requires $q\ne 1$.  If $q=1$
there is a neutral unit or a cancelled group because $|A|+k+q\ge 2$, and one
neutral pair is undone, so that $R$ becomes that pair (a same-parity
residual, handled in Step~4).  If one active child of free type remains and
$q=1$ or $q\ge 3$, a helper port labelled $-x$ absorbs it.  If $q=2$, two ports
labelled $a$ and $-x-a$ absorb it with a token of rank two.  If $q=0$ then
$|A|\ge 1$, so a neutral pair is undone and $v$ goes to Step~4 with the pair as
residual and the active child.  A remaining child of a joined type is never
absorbed by a helper (for the residue-one letter the helper would be $+x$,
which is the head of the grandchild).  Instead $v$ receives the context
$(h,\emptyset,p,c)$ with $c$ the joined type.

\paragraph{Step 3a (forced-antipode children).}
The exact census shows exactly two contexts whose every family contains the
antipode of its own head, the $(1;\text{one port})$ bottom and the two-input
cell with head chain of length two and residual $\{1\}$.  The first is folded
into its parent at Step~3 in every case, as the child type L11, so it never
becomes an owner, and the second is therefore the only forced-antipode child
that occurs.  Such a child $c$ is absorbed as follows.  If the parent is
stationary with at least two ports, the port block containing the head of $c$
absorbs it.  If the parent is an owner whose context has an input-in-token
family, that family absorbs it.  Alternatively, and this is what happens for
the contexts whose input-in-token search is infeasible, $c$ is joined into the
parent's cell as the child type A2, and the joint family is re-verified by the
checker as a whole.  The joined cell has a family that is not forced-antipode, so it
is an ordinary child of the next parent and the chain stops after one step.

A two-input owner can have two forced-antipode children, and it then joins one
of them and keeps the other pinned.  In the joint the pinned child contributes
one label whose output is that label itself, which is the role of a port, so
the joint is searched as the same context with one more port and child type
A2, and its family is required to carry the port in a token and to have no
label identically minus the port.  This is why we always join, and do not absorb,
when both are possible: a cell that is joined into its parent is then
a two-input cell with no forced-antipode children of its own, so the two
grandchildren that the joint pins are ordinary, and no cell ever needs more
than one token for an absorption.

The residue-one letter needs a word of its own, although it is not
forced-antipode.  With a direct parent edge and no port its own context has no
family with a free head, so it is folded into its parent as the child type L1.
With a subdivided parent edge, or with a port, it is an owner of its own.  However, when its
active child is forced-antipode the folding is not what we do: the
forced-antipode child is joined into the letter instead, which gives the
context $(h,\{1\},p,\mathrm{A2})$, and that context has a family that is not
forced-antipode, so the letter is an ordinary owner and its parent sees an
ordinary free child.  Therefore an L1 joint never carries a forced-antipode
grandchild.  (The opposite-parity bottom $(1,2)$ was joined in an earlier
version of the alphabet, although it has a family without the antipode of its
head and is now an ordinary owner.)

\paragraph{Step 4 (nonempty residual).}
$v$ is an owner with context $(h\bmod 6,R',p,c)$ where $R'$ is $R$ after the
following repairs: a portless residual $\{1\}$, $\{2\}$, $\{4\}$, $\{6\}$ or
$\{8\}$ with no active child retains one neutral unit, and a lone odd arm with
a head chain likewise retains one unit (its portless cell has only rays).  One
unit exists in each case because $v$ has at least two children.  The number
of recruited ports is $p=q$ when $q\le 2$, otherwise $p\in\{1,3\}$ with
$q-p\ne 1$, so that the unrecruited ports form a block of size zero or at
least two.  Every tuple produced in this way is in the list, itself
generated by the same rules.

\paragraph{Step 5 (root).}
The root has no head and its context is $(\mathrm{root},R',p,c)$ with the
same conventions, the root family containing the zero label.  Exact search
shows that a lone residual arm at the root needs two recruited ports, or none
for odd residue and one for even residue.  When the available ordinary
children do not permit this, one neutral pair is retained, and when the root
has a lone arm, no ports, no unit and $k\ge 2$ active children the exact root
table decides between cancelling all active heads and keeping one or two
active children.  A root whose only children are one residual arm and one
active child cannot occur, because branch vertices have three children.

Every owner is assigned exactly one context, every vertex of degree two lies
in exactly one arm, head chain or unit, and therefore the cells and units are
edge-disjoint and cover all subdivided edges.  \qed

The steps above are also the specification of our scheduler,
and the abstract enumeration described in Section~\ref{sec:decomp} is
our check that the list of contexts generated from these steps is complete:
the residual types, the port counts and the child types are finite, the
head-chain and arm lengths enter only through their residues modulo six
(with the raw lengths $1$ and $2$ distinguished), and every combination is
exercised by at least one fragment.


\section{Proof of the realization lemma}\label{app:realization}

\subsection*{Families, sinks and the choice rule}

Every family carries the properties (E1)--(E5) of Section~\ref{sec:cells},
which our checker verifies.  Property (E6) holds for a family of every context
that occurs except the \NumNotEsixContexts{} of Section~\ref{sec:cells}, and
for those the gate of Section~\ref{sec:verification} checks the finitely many
adjacent pairs individually.  For every owner we fix a family before any label
is chosen, preferring families with (E6).  With this choice no non-input label
of a parent is identically equal, as a form in the shared head, to a label of
its child: the parent's only input-only label is the antipode of its input, and
the child has no label equal to the antipode of its own head.

The tokens of a family are what (E2) leaves after the mirror pairs, but a
token may contain the input $x$ while the other token of the same family
contains its antipode $-x$ (which (E6) allows), or two tokens may otherwise
contain a label and its negative; such labels are already mirrored inside the
cell (the input $x$ is the head edge of the active child, $-x$ a label of the
cell) and must not be negated again.  We therefore work with \emph{units}:
take the multiset of the token labels of a family, remove every antipodal
pair of forms among them, and call each zero-sum triple that remains, or the
zero-sum set of four labels that remains when one pair straddled the two
tokens, a unit; what remains has size $6$ (two triples), $4$ (a \emph{quad}),
$3$ (a triple), $2$ or $0$ (nothing to place).  The first token of a
forced-antipode family, which contains the antipode of the head, keeps the
treatment below; every other unit is \emph{free}.  The zero-sum triples that
Section~\ref{sec:decomp} uses to cancel an odd number of active heads at a
vertex are free units as well.  They are not closed under negation, so their
negatives need sinks of their own; each of them consumes three active
children, all of which are owners, and triples at different vertices consume
disjoint sets of owners, which is what bounds their number in
Lemma~\ref{lem:packing}.  Free units are never
cancelled against one another along the root path.  Instead the negatives of
each free unit are placed in the ordinary part of the tree, at a \emph{sink},
and this is what closes the special support under negation.  There are two
kinds of sink.  A stationary vertex whose ordinary block has size $r$ takes
negated units directly as long as the remainder of $r$ after the units it
holds is zero or at least two, as Section~\ref{sec:dense} requires (a single
triple needs $r=3$ or $r\ge 5$, a single quad $r=4$ or $r\ge 6$); the units
are zero-sum, so the block remains zero-sum.  Otherwise units go to
\emph{gadgets} built from stationary vertices with two children: pairs of
triples to the \emph{absorbers} of Lemma~\ref{lem:absorber}, quads to the
four-vertex gadget of Lemma~\ref{lem:quad}.
Every cell chooses all of its parameters numerically, in post-order, so that
no symbolic parameter is ever live except in the two situations that already
had their own treatment, the forced-antipode children and the two-input rows.
This is a change of architecture with respect to an earlier version of this
appendix, in which tokens were cancelled in pairs along the root path by
solving $T_v=-T_A$ for the carrier of the resolving cell.  That step divided by
the determinant of a carrier block, and a chain of cells of the context
$\mathtt{h0\_3\_p0\_act}$ (a comb whose core vertices carry one arm each) has
a one-dimensional space of solutions, so that the labels either collide or
grow geometrically along the chain; Remark~\ref{rem:whysinks} records this.

\begin{lemma}[five-point absorber]\label{lem:absorber}
Let $v_0,v_1,v_2,v_3,v_4$ be vertices with $v_i$ the parent of $v_{i+1}$, none
of them an owner, a helper or a vertex of an arm, head chain or unit, such
that each $v_i$ with $i\le 3$ has a further child $w_i$ that is stationary,
$v_4$ has two such children $w_4,w_4'$, and the remaining children of every
$v_i$ number zero or at least two.  Let $T=(a,b,-a-b)$ and $U=(c,d,-c-d)$ be
integer zero-sum triples.  Give the edges the labels
\[
\begin{array}{lll}
e(v_0v_1)=a+b+d, & e(v_0w_0)=-a-b-d, \\
e(v_1v_2)=a,     & e(v_1w_1)=-c-d,   \\
e(v_2v_3)=b,     & e(v_2w_2)=d,      \\
e(v_3v_4)=2a+b-c,& e(v_3w_3)=c-2a-b, \\
e(v_4w_4)=c,     & e(v_4w_4')=-a-b,
\end{array}
\]
and give the remaining children of every $v_i$ a zero-sum block.  Then $v_0$
and $v_3$ are stationary, $v_1,v_2,v_4$ are non-stationary with
$s(v_1)=2a+b-c$, $s(v_2)=a+b+d$ and $s(v_4)=a$, which are the labels of the
parent edges of $v_4,v_1,v_2$, so that the identity of Lemma~\ref{lem:sigma}
holds on the five vertices with the parent edge of $v_0$ untouched; the ten
labels placed are $T\sqcup U\sqcup\{\pm(a+b+d)\}\sqcup\{\pm(2a+b-c)\}$; and as
linear forms in $(a,b,c,d)$ the ten labels are pairwise distinct and nonzero,
the only antipodal pairs among them being the two just named.  Distinctness of
the ten \emph{values} is then a finite list of proper conditions on
$(a,b,c,d)$, which the choice rule of Appendix~\ref{app:realization} adds to
$\Phi$ like any other; the lemma asserts it of the forms, not of every
numerical $T$ and $U$.
\end{lemma}

\begin{proof}
The sums are $s(v_0)=y+(a+b+d)-(a+b+d)=y$ for the parent label $y$ of $v_0$,
$s(v_1)=(a+b+d)+a-c-d=2a+b-c$, $s(v_2)=a+b+d$,
$s(v_3)=b+(2a+b-c)+(c-2a-b)=b$ and $s(v_4)=(2a+b-c)+c-a-b=a$.  A zero-sum
block at $v_i$ does not change $s(v_i)$, and a stationary child $w$ has
$s(w)=e(v_iw)$, so a hanging label is produced by its own lower endpoint
exactly as the label of a stationary child edge is.  The list of labels is
read off, and the distinctness of the ten forms and of their antipodal pairs
is a finite check, which we also ran by script.
\end{proof}

If two owners carry free tokens $-T$ and $-U$, the absorber places $T$ and $U$
and closes the special support under negation at the cost of two inverse
pairs, and it imposes no equation on either owner: the four coordinates
$a,b,c,d$ are whatever the two cells chose.  The three vertices $v_1,v_2,v_4$
join $W$.  On a vertical path of stationary vertices with two children each a
single token cannot be absorbed, since such a path places an even number of
labels, and three such vertices do not suffice for two tokens with independent
coordinates, because the three sums would have to permute the three path
labels and the parent label, which forces either an antipodal pair inside
$T\sqcup U$ or a dependence of the parent label on the tokens.  Four vertices
do not suffice either.  With path labels $g_1,g_2,g_3$ below the parent label
$y$ and hanging labels $p_0,p_1,p_2,p_3,p_3'$, the sum at the top vertex must be
$y$ (any other value would make $y$ depend on the tokens or place $-y$), so the
one inverse pair is $\{g_1,p_0\}$ and the six remaining labels are
$T\sqcup U$; the sums at the next two vertices cannot equal their own path
labels (that would force an antipodal pair inside $T\sqcup U$), so they are
$g_3$ and $g_1$ in this order, which gives $g_1+g_2+p_1=g_3$ and
$g_2+g_3+p_2=g_1$, hence $2g_2=-(p_1+p_2)$, and twice a label of $T\sqcup U$ is
never the sum of two others when the coordinates are independent.  Thus five
vertices is the minimum for two tokens on a binary path.

\begin{lemma}[quad gadget]\label{lem:quad}
Let $v_0,v_1,v_2,v_3$ be vertices with $v_i$ the parent of $v_{i+1}$,
eligible as in Lemma~\ref{lem:absorber}, each $v_i$ with $i\le 2$ having a
further stationary child $w_i$ and $v_3$ two such children $w_3,w_3'$.  Let
$Q=(q_1,q_2,q_3,q_4)$ be an integer zero-sum quadruple and $f$ an integer.
Give the edges the labels
\[
\begin{array}{lll}
e(v_0v_1)=f-q_1-q_2, & e(v_0w_0)=-(f-q_1-q_2), \\
e(v_1v_2)=q_1,       & e(v_1w_1)=q_2,          \\
e(v_2v_3)=f,         & e(v_2w_2)=-f,           \\
e(v_3w_3)=q_3,       & e(v_3w_3')=q_4.
\end{array}
\]
Then $v_0$ is stationary, $v_1,v_2,v_3$ are non-stationary with sums
$f$, $q_1$ and $f-q_1-q_2$, the labels of the parent edges of $v_3,v_2,v_1$
respectively, so that the identity of Lemma~\ref{lem:sigma} holds on the four
vertices with the parent edge of $v_0$ untouched; the eight labels placed are
$Q\sqcup\{\pm f\}\sqcup\{\pm(f-q_1-q_2)\}$, and they are pairwise distinct and
nonzero as forms in $(q_1,q_2,q_3,f)$.
\end{lemma}

\begin{proof}
$s(v_1)=(f-q_1-q_2)+q_1+q_2=f$, $s(v_2)=q_1+f-f=q_1$ and
$s(v_3)=f+q_3+q_4=f-q_1-q_2$; the rest is as in Lemma~\ref{lem:absorber}.
\end{proof}

The parameter $f$ is ours, chosen when the gadget is filled, so the
conditions that its four labels $\pm f$, $\pm(f-q_1-q_2)$ avoid every placed
value are proper affine conditions in $f$, and a quad gadget occupies the top
four vertices of any five-vertex path of the kind Lemma~\ref{lem:packing}
provides.

\begin{remark}[a quad splits into two triples]\label{rem:quadsplit}
The negatives $Q=(q_1,q_2,q_3,q_4)$ of a free quad need not go to a gadget or
to a block that can hold four labels.  They can be placed as the two triples
$\{q_1,q_2,-q_1-q_2\}$ and $\{q_3,q_4,-q_3-q_4\}$ at two ordinary triple
sinks.  Both sets are zero-sum, so both blocks stay zero-sum, and because $Q$
is zero-sum the two labels introduced are $-q_1-q_2$ and its negative, an
inverse pair; the special support therefore stays closed under negation and
the ledger gains nothing that needs a sink of its own.  Each new label is a
sum of two coordinates of $Q$, so the magnitude bound at the end of this
appendix covers it, and each carries one further genericity condition, of the
same kind as the conditions on $f$ above.  We therefore count every quad as
\emph{two} triples in Lemma~\ref{lem:packing}.  This is what makes
$\mathrm{cap}$, which counts triples, the right capacity for an arbitrary
mixture of triples and quads: a block of size three or five has capacity one
but cannot hold a quad at all, so a capacity counted in units rather than in
triples would not do.
\end{remark}

\begin{lemma}[sinks exist]\label{lem:packing}
Let $t$ be the number of free units, a quad counted as two triples
(Remark~\ref{rem:quadsplit}).  If $n\ge 10^6(D+1)$ then the
direct capacities of the ordinary blocks, $\sum_u\mathrm{cap}(r_u)$ with
$\mathrm{cap}(r)=\max\{k: r-3k\in\{0\}\cup[2,\infty)\}$, together with the
number of pairwise vertex-disjoint vertical paths of five stationary vertices
whose ordinary blocks have size two or four, is at least $t$.  (A path takes two triples, or one quad, which counts
for two; a block of size $r$ takes a quad in place of two triples whenever
$r-4\in\{0\}\cup[2,\infty)$, and otherwise the quad is split as in
Remark~\ref{rem:quadsplit}.  Counting a quad as two triples on either route
is therefore conservative.)
\end{lemma}

\begin{proof}
By Lemma~\ref{lem:charge} there are at most $2q\le 2D$ owners, each with at
most two units of its own, and a unit counts for at most two triples, so the
units of the families contribute at most $8D$.  A cancelling triple of
Section~\ref{sec:decomp} consumes three owners, and triples at different
vertices consume disjoint sets of them, so there are at most $2D/3$ of those.
Hence $t\le 8D+2D/3\le 9D$; helpers are at most one per owner.  The skeleton, by which we
mean the vertices of $W$, the arms with their leaves, the units, the head
chains, the owners, the helpers, the helper ports (at most two per helper)
and the recruited ports (at most three per owner), has at most $16D+1$
vertices.  Suppose the direct capacities sum to less than
$t$.  Then fewer than $t$ blocks have positive capacity, and a block of
capacity $k\ge 1$ has size at most $3k+4$ (otherwise $k+1$ would be
admissible), so these blocks have fewer than $7t\le 63D$ children in total.
Every vertex outside the skeleton is a child in exactly one ordinary block,
and blocks of size one do not exist, so all but $7t$ of the $n-16D-1$
non-skeleton vertices are children in blocks of size two or four.  Call a
vertex \emph{usable} if it is neither an owner, nor a helper, nor on an arm,
head chain or unit, and its block has size two or four; let $U$ be the set of
usable vertices and $m=|U|$.  The parent of a child in a block of size two or
four is usable or in the skeleton, and has at most four such children, so
$n-16D-1-63D\le 4(m+16D+1)$, i.e.\ $m\ge n/8$ once $n\ge 286D+10$.

Consider the forest $F$ induced on $U$.  A vertex of $U$ has at most four
children in $U$, and the roots of $F$ are children of skeleton vertices or of
blocks with positive capacity, so $F$ has at most $4(16D+1)+63D=127D+4$
components.  Give
each vertex of $F$ its height, leaves of $F$ having height $0$.  A vertex of
height at most $4$ lies either in a component of height at most $4$, which
has at most $1+4+16+64+256=341$ vertices, or below a child of height at most
$4$ of a vertex of height at least $5$, and every vertex of height at least
$5$ has at most four such children, each with at most $341$ vertices below
it.  Writing $H$ for the number of vertices of height at least $5$, we get
$m\le 1365H+341(127D+4)$, hence $H\ge m/2730$ once $m\ge 86614D+2728$, which
holds when $n\ge 7.0\cdot 10^5(D+1)$.  For a vertex $u$ of height $h\ge 5$ with
$h\equiv 0\pmod 5$ follow the path from $u$ downward through children of
maximal height; it visits heights $h,h-1,\dots,h-4\ge 1$, five usable
vertices; the first four have a child on the path and therefore, since their
blocks have size two or four, a further child on which a gadget label can
hang, and the last has height at least one, so its block of size two or four
consists of children off the path, two of which take the last two labels.  A
hanging child may itself be the top vertex of another gadget; it is then
stationary all the same, which is all that is required of it.  Two such paths are disjoint, because
from incomparable vertices they lie in disjoint subtrees, and from comparable
vertices $u$ above $u'$ the heights satisfy $h(u)\ge h(u')+5$, so the height
ranges do not meet.  Thus there are at least $H/5\ge m/13650\ge n/109200$
disjoint paths, which is at least $t\le 9D$ once $n\ge 982800D$.  The three
thresholds are implied by $n\ge 10^6(D+1)$.
\end{proof}

The constant is crude, and it is irrelevant beside $A$ below.  In our
experiments the direct capacity alone covers the tokens on random sparse trees
with $n\ge 20(D+1)$, e.g.\ by a factor of at least twelve at $D\approx n/25$,
and on a comb attached to a complete binary tree of depth $d$, where only
absorbers can help, the greedy packing finds $2^{d-4}$ disjoint paths.  Blocks at owners (their unrecruited ports) and at helpers (their
further ports) count toward the direct capacity as well.  A usable vertex with
a block of size four gives its two further children an inverse pair, which
does not change its sum.

\begin{lemma}[compatibility]\label{lem:compat}
Send every free quad to a quad sink and pair the free triples arbitrarily,
with at most one left over; a dependent pair of triples (two units of one
cell, or of a cell and its active child) is given a role order for which none
of the four absorber forms $\pm(a+b+d)$, $\pm(2a+b-c)$ is identically equal,
over the joint parameters of the two cells, to a label of either cell, to the
negative of one, to another absorber form or to zero.  Such an order exists
for every pair of catalogue families.  Then, when the owners are processed in
post-order with all parameters numeric, every condition that has to be
avoided, including the distinctness of the negated units placed at the sinks
from every other label, is a proper affine condition in the coordinate to
which it is charged, and Lemma~\ref{lem:charge} applies.
\end{lemma}

\begin{proof}
The labels of a family are homogeneous linear forms, pairwise distinct and
nonzero by (E3).  We first dispose of the negated units themselves.  Let $u$
be a label of a free unit of the cell $v$ and consider the value $-u$ placed
at the sink.  It is not identically a label of $v$: the units were formed by
removing every antipodal pair among the token labels, and a label outside the
tokens is a mirror-pair member, whose negative is its partner and not $-u$.
It is not identically a label of the active child of $v$, whose labels are
forms in that child's parameters: the only parameter shared is the child's
head $x$, and when both cells satisfy (E6) the child has no non-head label
that is a multiple of $x$, while $u=-x$ would mean that $x$ and $-x$ both
occur among the tokens of $v$ and were removed.  It is not identically a
label of the parent of $v$, which shares only the head $y$ of $v$: under (E6)
the parent's labels depending on $y$ alone are $\pm y$, and $u=-y$ occurs
only in the first token of a forced-antipode family, which is not free.  For
the finitely many pairs in which a cell violates the second clause of (E6)
(a unit label that is a multiple of the input by a factor other than $\pm1$,
of which the catalogue has $30$ in $28$ contexts, with factors $\pm2$, $\pm3$,
$\tfrac12$, $\tfrac32$; or a label that is such a multiple of the head, of
which there are $478$ in $166$ contexts) the identity is checked directly, as
part of the gate of Section~\ref{sec:verification}, and the choice rule keeps
the colliding combinations apart.  Every other label of the tree is numeric
when the sink is filled, so the remaining requirements on $-u$ are proper
conditions on the parameters of $v$.  The four labels of a quad gadget are affine in the fresh parameter $f$
with slope $\pm 1$, so their conditions are proper in $f$.  Consider now an
absorber whose first source is already numeric
and whose second source is the present cell with token $(c,d,-c-d)$.  Each
absorber form is a numeric constant plus one of $d,-c,-d,c$ (or, with the
roles exchanged, plus $a+b$, $2a+b$ and their negatives), so it has a nonzero
coefficient in some parameter of the present cell.  Its collision with a
placed numeric label is therefore a proper condition.  Its collision with a
label $\ell$ of the present cell reads $d-\ell=\mathrm{const}$, which is proper
unless $d=\ell$ identically, and (E3) excludes that; collisions among the four
absorber forms and the two tokens read $c\pm d=\mathrm{const}$, $c=\mathrm{const}$
or $d=\mathrm{const}$, again proper because $c,d,-c-d$ are distinct nonzero
labels.  If the two sources are independent this is all.  If they are
dependent, the difference of two forms is a form over the joint parameters; it
is proper for the cell processed second if it involves one of that cell's
parameters, it is a proper condition on the cell processed first if it
involves only that cell's parameters (and the pairing is known before any
label is chosen, so the condition is recorded in time), and only an identical
zero would be fatal.  The finite check over every pair of a family with one
input and a first token and a family with a head and a first token, over every
entry of every menu, $9{,}637{,}425$ ordered pairs in all, finds a
collision-free role order for each of them, in fact at least six of the
thirty-six; the same check for two-input parents,
second tokens and the two tokens of one cell is described in
Section~\ref{sec:verification}.
\end{proof}

Without the finite check one would pair independent tokens only, which is
possible for $t\ge 10$ because the dependency graph has maximum degree four
and its complement then has minimum degree at least $t/2$; the check makes
this unnecessary.

\begin{remark}[why the tokens are not cancelled along the root path]\label{rem:whysinks}
An earlier version of this appendix cancelled a free token against the first
token on the path to the root by solving $T_v=-T_A$ for the carrier of the
resolving cell.  Three things go wrong with this, and we record them because
each was found only after the argument had been written out.  First, when
the pending cell is the active child of the resolving cell, the child's head
is the parent's input and appears on both sides of the equation, so the
displayed solution is not a definition; the correct statement is two
equations in three unknowns, whose relevant minors are mixed ones and not the
determinant of a carrier block.  Second, the division by that determinant was
measured over the coordinates $(y,x)$ with the family's denominator kept
apart; measured on the lattice of parameter values at which the family's
labels are integers, the token of $\mathtt{h0\_3\_p0\_act}$ is a bijection
onto the integer zero-sum triples and nothing is divided, whereas the tokens
of $\mathtt{h1\_2\_p0\_act}$ and $\mathtt{h4\_1\text{-}3\text{-}3\text{-}3\_p0\_act}$
have image index two and three.  Third, and decisively, the four contexts
whose tokens depend on the head and the input alone have one free parameter
per cell, so a chain of $m$ of them, which our scheduler produces on a comb
with one arm of length three at every core vertex, has a one-dimensional space
of solutions once the tokens are paired, whatever the pairing: the heads
either collide or grow by a factor of at least $5/2$ per pair, and our
proof-faithful constructor records the failure on such a comb attached to a
dense part.  A two-level joint of two such cells does have a tokenless family
with a free head (we verified this for all nine ordered pairs among the three
small contexts), which would cut the chains as Lemma~\ref{lem:reset} cuts the
two-input chains.  However, sending the tokens to sinks removes the pairing
altogether, and with it the division, the chains and the live sets, so we use
the sinks.
\end{remark}

\subsection*{State}

Owners are processed in post-order.  Before any label is chosen the free
tokens are paired and every token or pair is assigned a sink by
Lemma~\ref{lem:packing}, and every dependent pair is given a collision-free
role order by Lemma~\ref{lem:compat}.  The state consists of the set $\Placed$
of numeric special labels, the set of \emph{reserved} values (antipodes of
heads that a parent will place later), the set $\mathrm{Open}\subseteq\Placed$
of placed labels whose antipode is not yet placed and is designated (the head
of an active child, the three labels of a free token whose sink is not yet
filled), and the set $\Phi$ of recorded forbidden conditions.  The only
symbolic labels are those of a forced-antipode child before its parent is
processed, of a two-input row of residual $\{1\}$ before its reset, and the
four labels of an absorber whose second source is the present cell; each is
an affine form in the parameters of at most two adjacent cells.

A numeric value $v$ is \emph{admissible} for a label if $v\ne 0$, $v$ is not
placed or reserved, and $-v$ is not placed or reserved, unless $-v$ is open and
$v$ is its designated antipode (the forced $-x$ of an absorbing family, a
helper, or a coordinate of a negated token when a sink is filled).

The invariant is: (a) numeric labels are pairwise distinct and nonzero, and no
two are antipodal except designated pairs; (b) every symbolic label is a
non-constant affine form in the parameters of the present cell or of its
active child, and no two symbolic labels are identical or antipodal as forms;
(c) $\Phi$ contains every condition under which a symbolic label would collide
with a placed value or with another symbolic label; (d)
$|\Placed|+|\mathrm{Reserved}|\le|S|+\#\text{cells}$ and $|\Phi|\le 3W(2|\Placed|+3W)$.

\subsection*{Processing an owner}

Let $v$ be the next owner with family $F$.  Its input is the head label of the
remaining active child, which is numeric, since every cell chooses its head
numerically; the one exception is a child of the two-input residual-$\{1\}$
type, treated below.

\emph{No token.}  Choose the head and the internal parameters numerically,
one coordinate at a time, in $M\Z\cap[-B,B]$, in a fixed order of the
parameters.  Property (E3) says that each label form is nonzero and that no two
of them are equal, so each label form, and each difference of two label forms,
has a well-defined \emph{last} nonzero coordinate in that order.  We charge
every admissibility requirement and every reservation to the step at which the
last nonzero coordinate of the form concerned is chosen; at that step the form
is affine in the coordinate being chosen with nonzero slope, whereas at every
later step it is already numeric.  Each requirement therefore excludes at most
one value, at exactly one step, and the reservation is
$\{q y: q\in\Phi(F_{\mathrm{parent}})\}$.  We note that (E3) alone would not
justify this coordinate by coordinate, since a label may well have zero
coefficient in several parameters; it is the charging to the last nonzero
coordinate that makes the count valid.  The bound
$B=M\bigl(2W(|\Placed|+|\mathrm{Reserved}|+1)+W^2+1\bigr)$ therefore suffices.

\emph{Free token.}  Choose all parameters numerically exactly as in the
previous case, with two additions to the forbidden set.  If the token's sink
is an absorber whose other source is already placed, the four absorber forms
of Lemma~\ref{lem:absorber}, with the other token's coordinates numeric and
this token's coordinates as forms in the present parameters, are added to
$\Phi$, together with their negatives; by Lemma~\ref{lem:compat} each is a
proper affine condition, and there are at most eight of them.  If the two
sources are dependent and the present cell is the one processed first, the
conditions of Lemma~\ref{lem:compat} that involve only the present cell's
parameters are added as well.  When the last source of a sink is placed, the
sink is filled: a direct sink receives the negated triple as three prescribed
members of its block, and an absorber receives its ten labels, of which the
six of $T\sqcup U$ are the designated antipodes of the two tokens and the four
others are fresh; all become placed.  The three labels of a token are open
from the moment they are placed until the sink is filled.

\emph{Two-input rows.}  Treated identically with both inputs.  The rows we
call rigid need a word here, because their heads are written in
Appendix~\ref{app:contexts} as fixed forms in their inputs, and so are their
other labels.  Our constructor carries them as templates outside the
catalogue: the residue-two ray as $(-3t;\,-t,\,-2t,\,t;\,2t)$ with input
$x=2t$ and head $-3x/2$, and the direct-edge two-input rows with residual
$\{1\}$ and $\{2\}$ as $(-x_1-x_2;\,-x_2,\,-x_1;\,x_1,x_2)$ and as the
$\tfrac12$-form of Appendix~\ref{app:contexts}.  (Solving $P'=O$ for the
residue-two topology does give a two-dimensional space containing the
template, but its second direction does not meet the signature of
Section~\ref{sec:cells}, so we do not use it, and the template's head is not a
free parameter.)  A rigid row has no free unit, and all of its labels are
determined by its inputs; every requirement on such a label is charged to the
step at which the last of those inputs is chosen, exactly as the
coordinate-by-coordinate rule charges a label to its last coordinate, and the
labels of a rigid row are reserved from that step on.  One more thing has to
be checked for the residue-two ray, whose non-head labels are
$x/3,2x/3,-x/3,-2x/3$ in terms of the head $x$ that its parent receives as
input: a parent family with a label that is one of these multiples of its
input would collide identically.  The choice rule therefore excludes such
families above a residue-two ray; Section~\ref{sec:verification} reports the
count of families and contexts concerned.  For the row with residual $\{1\}$
an exhaustive enumeration of the bijections between the five labels and the
five outputs of that topology leaves exactly one nondegenerate family, and
its head is $-x_1-x_2$.  That row has no token, so a chain of such rows
carries a symbolic head upward for as many levels as the chain is long.  Our scheduler does produce arbitrarily long such
chains, on a comb whose every core vertex has two active children and one arm
of length one, so the length is $\Theta(D)$ in the worst case and not $O(1)$.
What bounds it is Lemma~\ref{lem:reset}.

\begin{lemma}[the two-input reset]\label{lem:reset}
Two consecutive residual-$\{1\}$ direct-edge two-input rows, joined into a
single cell, have a family whose head is an independent parameter.
\end{lemma}

\begin{proof}
The joint has nine labels, the heads $g_0$ and $g_1$ of the two rows, the two
labels of each arm, the pinned input $s_0$ of the upper row and the two pinned
inputs $x_1,x_2$ of the lower one, and nine outputs.  Enumerating every
bijection between them and solving the resulting linear system gives $14$
families whose labels are distinct and nonzero at a generic point, and for $12$
of them the space of solutions contains a direction along which $g_0$ moves
while every input is held fixed.  One such family, at an integer point, is
$$(g_0,g_1)=(14,-16),\quad
  (r_0a_0x_0,r_0a_0x_1)=(13,-29),\quad
  (r_1a_0x_0,r_1a_0x_1)=(-9,22),$$
$$s_0=3,\quad (x_1,x_2)=(-3,19),$$
whose solution space has dimension five and whose head-free direction is $g_0$
itself, so $g_0$ may be moved on its own.  We verified $P=O$ and the
distinctness of the nine labels at that point directly.  (This family is the
one with the fewest moving labels once the inputs are pinned; the constructor
uses a different one of the twelve, see below.)
\end{proof}

A chain of such rows is therefore cut into joints of two levels, each of which
has a free head, and a chain of odd length leaves one single row at its top
whose parent is not one of these rows.  The symbolic labels are confined to at
most two cells and one further row, where a cell may now be a joint of width
nine, so the bound is a constant and the boxes and the magnitudes are again
linear in $D+1$.  Of the $12$ head-free families, seven have the further
property that all six labels placed by the joint still move once its three
inputs are pinned, and our constructor uses one of these; for the single row
that an odd chain leaves at its top, whose labels are determined by its inputs,
the requirements are charged to the step at which the last of those inputs is
chosen, exactly as in the coordinate-by-coordinate rule above.  Our scheduler
reports these joints (field \texttt{reset\_joints}) and the constructor of
Section~\ref{sec:verification} realizes them as single cells.

\emph{Stationary vertices.}  A helper places $-x$, a two-port absorption
places a fresh numeric $a$ and $-x-a$, and neutral units are instantiated with
fresh numeric parameters.  A stationary vertex that is a sink is filled as
described under \emph{Free token}, and its remaining children, if any, are
left to the completion as a zero-sum block of size at least two.

\emph{Two-token cells.}  The units of the cell (two triples, or one quad when
the tokens share an antipodal pair) are sent to sinks like any other; they are
dependent on each other and on the units of the parent and of the active
children, which Lemma~\ref{lem:compat} allows.

\emph{Forced-antipode children.}  The token $\{-x,a,x-a\}$ of a
forced-antipode child contains the antipode of the child's own head edge $x$,
whose negative is that edge itself, so it cannot be sent to a sink.  It stays
symbolic with its carrier live until the parent is processed.  The parent's
port block or input-in-token family supplies the token $\{x,b,-x-b\}$, and the
single equation $a=-b$ is solved for the child's carrier (its token map has
rank two on the carrier, so the equation is non-degenerate, and it involves no
division since $b$ is a single label).  The second token of a forced-antipode
family, when there is one (the families of the one forced-antipode context
that occurs all carry two, the second being $\{x_1,x_2,-x_1-x_2\}$ on the two
inputs), is a free unit and goes to a sink, as is a second token of an
input-in-token family of the parent.

\emph{Root.}  The root family contains the zero label; its free tokens go to
sinks like any other.  If the number of free tokens is odd, one token is left
over by Lemma~\ref{lem:compat}; its negatives are handed to an odd ordinary
block of size at least three, which exists because the special support then
has odd size, so the ordinary block sizes have odd sum, and no block has size
one.

\subsection*{Verification of the invariant and of the bounds}

(a) follows from admissibility; (b) from (E3) and the choice rule for the
forced-antipode and two-input cases and from Lemma~\ref{lem:compat} for the
absorber forms; (c) by construction; (d) because every owner adds at most $W$
placed labels and at most $|\Phi(F)|\le\NumPhiMax$ reservations, the bound
$\NumPhiMax$ being computed from the catalogue as in Lemma~\ref{lem:charge},
every absorber adds four placed labels and a direct sink none, and every
absorber adds at most eight conditions to $\Phi$.  Every box has side
$O(|\Placed|+|\mathrm{Reserved}|+|\Phi|)=O(|S|)$, every label is a
bounded-coefficient combination of at most seven parameters, and
$|S|\le 12D+3+2t\le 30D+3$ by Lemma~\ref{lem:charge} and $t\le 9D$, so all
magnitudes are at most $A(D+1)$ for an absolute $A$.  Explicitly, with
$W=\NumMaxWidth$ and $M=\NumLatticeM$, and with
$|\Placed|+|\mathrm{Reserved}|$ together with the absorber conditions at most
$(30+2+2\NumPhiMax+8)(D+1)=\NumPlacedCoefSink(D+1)$, the coordinate boxes have
side at most
$M\bigl(2W(\NumPlacedCoefSink(D+1)+1)+W^2+1\bigr)=\NumSinkBoxLin(D+1)+\NumSinkBoxCon$.
A label of a cell is a sum of at most $\NumParams$ terms with coefficients at
most $\NumCoefBound$ over a denominator that only decreases the magnitude, so
it has magnitude at most
$\NumParams\cdot\NumCoefBound\cdot(\NumSinkBoxLin(D+1)+\NumSinkBoxCon)
=\NumSinkLabLin(D+1)+\NumSinkLabCon$, and an absorber label is a sum of at most
four token coordinates, so every special label has magnitude at most
$4(\NumSinkLabLin+\NumSinkLabCon)(D+1)\le\AConst(D+1)$.  This is
the constant $A$ of Lemma~\ref{lem:realization}; no lattice or coefficient
enlargement occurs, since nothing is ever divided.  At the end every special
label other than a possible single leftover triple has its antipode placed,
so the special support is closed under negation up to that triple.  \qed


\bibliographystyle{plainnat}
\bibliography{refs}

\end{document}
